\documentclass[../main.tex]{subfiles}
\begin{document}

\chapter{スケジューリング}
\lessoninfo{
  video={P3L1},
  textbook={OSTEP \cite{ostep} ch.~7--10；Silberschatz ら \cite{silberschatz2018} ch.~5；Fedorova ら \cite{fedorova2004}},
  keywords={ランキュー，FCFS，SJF，優先度逆転，ラウンドロビン，タイムスライス，MLFQ，O(1) スケジューラ，CFS，キャッシュアフィニティ，NUMA，同時マルチスレッディング，ハードウェア性能カウンタ},
}

第2章では，レディキューから次に実行するプロセスを選ぶ CPU スケジューラを
導入し，タイムスライスの長さがスケジューリングのオーバーヘッドと，実行
可能なプロセスが待つ時間とを釣り合わせることを示した．この章では，その
選択がどのように行われるかを調べる．単純な実行完了型の方針から，
プリエンプティブなもの，優先度に基づくもの，タイムシェアリングを行うものに
至るスケジューリングアルゴリズム，それらを実装するランキューのデータ構造，
そして Linux カーネルの2つのスケジューラを扱う．続いて，マルチプロセッサ
とハイパースレッディング対応のプロセッサが提起する問題を，Fedorova ら
\cite{fedorova2004} の論文に沿って論じる．アルゴリズムの比較には第6章の
性能指標を用いる．

\section{CPU スケジューリングの概観}
\label{sec:l07-overview}

CPU スケジューラは，プロセスと，その中のスレッドが，共有された CPU に
いつどのようにアクセスするかを決定する．\source{P3L1，視覚的な比喩，
スケジューリングの概観；OSTEP ch.~7；Silberschatz ら ch.~5．}ユーザレベル
のプロセスやスレッドを実行するタスクも，カーネル自身のカーネルレベル
スレッドも，ともにスケジュールする．第5章の Linux の用語にならい，この章
ではそのようなスケジュール可能な実体をすべて\emph{タスク}と呼ぶ．

最も単純なスケジューラでさえ1つの目標を体現している．3つの初歩的な戦略が，
ありうる目標の幅を示す．
\begin{itemize}
  \item タスクを到着した順に直ちにディスパッチする．スケジューリングの
    判断は自明であり，そのオーバーヘッドは小さい．
  \item 単純で短いタスクを先にディスパッチする．これは早い時期に最も多く
    のタスクを完了させ，したがってワークロード全体より短いどの区間に
    ついても，単位時間当たりに完了するタスク数であるスループットを最大化
    する．
  \item 複雑で長いタスクを先にディスパッチする．これは資源を占有し続け，
    CPU，デバイス，メモリの使用率を最大化する．
\end{itemize}

スケジューラは，CPU を別のタスクに与えなければならない可能性があるときに
必ず動作する．
\begin{itemize}
  \item CPU がアイドルになったとき．例えば実行中のタスクが I/O を待って
    ブロックした場合である．
  \item 新たなタスクが実行可能になったとき．例えば新しく生成されたタスク
    や，I/O 操作が完了したタスクである．スケジューラは，その新しいタスク
    が実行中のタスクより重要であって，それに割り込むべきかどうかを判断
    しなければならない．
  \item 実行中のタスクのタイムスライスが満了したとき．\margin{満了に
    気付くのはタイマのティックである．Linux カーネルは \texttt{CONFIG\_HZ}
    を 100，250，1000 のいずれかで構築され，ティックは 10，4，
    \SI{1}{\milli\second} ごとに入る．}
\end{itemize}
そしてスケジューラは実行可能なタスクの1つを選んでディスパッチする．選ばれ
たタスクへのコンテキストスイッチを行い，ユーザモードに入り，プログラム
カウンタをそのタスクの次の命令に設定する．タスクはそこから実行を続ける．

どのタスクが選ばれるかはスケジューリング方針，すなわちアルゴリズムが
決める．選択がどのように実行されるかは，実行可能なタスクを保持する
データ構造に依存する．

\begin{definition}[ランキュー]
  スケジューラが実行可能なタスクを保持しておくデータ構造を，その
  \term[らんきゅー]{ランキュー}[runqueue]という．
\end{definition}

第2章のレディキューはそのようなランキューの1つである．名前に反して，
ランキューは FIFO キューである必要はない．多くのスケジューラは木や複数の
キューを用いる．ランキューとスケジューリングアルゴリズムは密接に結び付いて
いる．アルゴリズムを効率的に実装できるのは，その選択を安価にするランキュー
がある場合に限られ，ランキューの側も，それが支えるアルゴリズムに合わせて
設計される．

\needspace{7\baselineskip}
\section{実行完了型スケジューリング}
\label{sec:l07-rtc}

最初に考えるスケジューラは，タスクから CPU を決して取り上げない．
\source{P3L1，実行完了型スケジューリング，SJF の性能；OSTEP ch.~7；
Silberschatz ら ch.~5．}これらを単純化した仮定の下で解析する．一群の
タスクをスケジュールするものとし，各タスクの実行時間はあらかじめ分かって
おり，プリエンプションはなく，CPU は1つであるとする．アルゴリズムは，
スループット，平均完了時間，平均待ち時間，CPU 使用率（第6章）によって
比較する．タスクの完了時間はその到着から完了までで測り，待ち時間はその
到着から実行が始まるまでで測る．

\begin{definition}[実行完了型スケジューリング]
  \term[じっこうかんりょうがたすけじゅーりんぐ]{実行完了型スケジュー%
  リング}[run-to-completion scheduling]では，いったんディスパッチされた
  タスクは完了するまで CPU 上で実行され，プリエンプトされることはない．
\end{definition}

\begin{definition}[FCFS]
  \term[せんちゃくじゅんすけじゅーりんぐ]{先着順スケジューリング}%
  [first-come first-served scheduling]%
  \index{FCFS|see{先着順スケジューリング}}（FCFS）に従うスケジューラは，
  タスクを到着した順に実行する．
\end{definition}

FCFS のランキューは FIFO キューである．到着したタスクは末尾に追加され，
スケジューラは先頭のタスクを取り出す．どちらの操作も定数時間で済む．

\begin{definition}[SJF]
  \term[さいたんじょぶゆうせん]{最短ジョブ優先}[shortest job first]%
  \index{SJF|see{最短ジョブ優先}}（SJF）という方針に従うスケジューラは，
  タスクを実行時間の短い順に実行する．
\end{definition}

SJF のランキューは実行時間で順序付けられていなければならない．順序付き
キューを使えばスケジューラは先頭のタスクを定数時間で取り出せるが，タスクの
挿入にはその位置の探索が必要で，キューの長さに比例する時間がかかる．実行
時間で順序付けた平衡木ならタスクの挿入は対数時間で済み，スケジューラは
最左の節点を取り出す．

\begin{figure}[htbp]
  \centering
  % FCFS, SJF and round robin (timeslice 1 s) for three tasks that arrive
% at time 0 in the order T1, T2, T3 with execution times 5 s, 2 s, 1 s.
% Usage: % FCFS, SJF and round robin (timeslice 1 s) for three tasks that arrive
% at time 0 in the order T1, T2, T3 with execution times 5 s, 2 s, 1 s.
% Usage: % FCFS, SJF and round robin (timeslice 1 s) for three tasks that arrive
% at time 0 in the order T1, T2, T3 with execution times 5 s, 2 s, 1 s.
% Usage: \input{figures/l07-gantt}   (width ~106 mm)
\begin{tikzpicture}[x=10mm, y=1mm, font=\footnotesize\sffamily,
  t1/.style={draw, fill=ln-accent!25},
  t2/.style={draw, fill=ln-box},
  t3/.style={draw, fill=white},
  lab/.style={anchor=east, align=right}]
  % FCFS
  \node[lab] at (-0.2,25) {FCFS};
  \draw[t1] (0,22) rectangle (5,28); \node at (2.5,25) {T1};
  \draw[t2] (5,22) rectangle (7,28); \node at (6,25) {T2};
  \draw[t3] (7,22) rectangle (8,28); \node at (7.5,25) {T3};
  % SJF
  \node[lab] at (-0.2,14) {SJF};
  \draw[t3] (0,11) rectangle (1,17); \node at (0.5,14) {T3};
  \draw[t2] (1,11) rectangle (3,17); \node at (2,14) {T2};
  \draw[t1] (3,11) rectangle (8,17); \node at (5.5,14) {T1};
  % round robin, timeslice 1 s
  \node[lab] at (-0.2,3) {\iflangja{RR（1\,s）}{RR (1\,s)}};
  \draw[t1] (0,0) rectangle (1,6); \node at (0.5,3) {T1};
  \draw[t2] (1,0) rectangle (2,6); \node at (1.5,3) {T2};
  \draw[t3] (2,0) rectangle (3,6); \node at (2.5,3) {T3};
  \draw[t1] (3,0) rectangle (4,6); \node at (3.5,3) {T1};
  \draw[t2] (4,0) rectangle (5,6); \node at (4.5,3) {T2};
  \draw[t1] (5,0) rectangle (8,6); \node at (6.5,3) {T1};
  \draw[densely dotted] (6,0) -- (6,6) (7,0) -- (7,6);
  % time axis
  \draw[->] (0,-3) -- (8.6,-3);
  \node[anchor=west] at (8.6,-3) {$t$\,(s)};
  \foreach \x in {0,...,8} {
    \draw (\x,-3) -- (\x,-4.2);
    \node[below, font=\scriptsize\sffamily] at (\x,-4.2) {\x};
  }
\end{tikzpicture}
   (width ~106 mm)
\begin{tikzpicture}[x=10mm, y=1mm, font=\footnotesize\sffamily,
  t1/.style={draw, fill=ln-accent!25},
  t2/.style={draw, fill=ln-box},
  t3/.style={draw, fill=white},
  lab/.style={anchor=east, align=right}]
  % FCFS
  \node[lab] at (-0.2,25) {FCFS};
  \draw[t1] (0,22) rectangle (5,28); \node at (2.5,25) {T1};
  \draw[t2] (5,22) rectangle (7,28); \node at (6,25) {T2};
  \draw[t3] (7,22) rectangle (8,28); \node at (7.5,25) {T3};
  % SJF
  \node[lab] at (-0.2,14) {SJF};
  \draw[t3] (0,11) rectangle (1,17); \node at (0.5,14) {T3};
  \draw[t2] (1,11) rectangle (3,17); \node at (2,14) {T2};
  \draw[t1] (3,11) rectangle (8,17); \node at (5.5,14) {T1};
  % round robin, timeslice 1 s
  \node[lab] at (-0.2,3) {\iflangja{RR（1\,s）}{RR (1\,s)}};
  \draw[t1] (0,0) rectangle (1,6); \node at (0.5,3) {T1};
  \draw[t2] (1,0) rectangle (2,6); \node at (1.5,3) {T2};
  \draw[t3] (2,0) rectangle (3,6); \node at (2.5,3) {T3};
  \draw[t1] (3,0) rectangle (4,6); \node at (3.5,3) {T1};
  \draw[t2] (4,0) rectangle (5,6); \node at (4.5,3) {T2};
  \draw[t1] (5,0) rectangle (8,6); \node at (6.5,3) {T1};
  \draw[densely dotted] (6,0) -- (6,6) (7,0) -- (7,6);
  % time axis
  \draw[->] (0,-3) -- (8.6,-3);
  \node[anchor=west] at (8.6,-3) {$t$\,(s)};
  \foreach \x in {0,...,8} {
    \draw (\x,-3) -- (\x,-4.2);
    \node[below, font=\scriptsize\sffamily] at (\x,-4.2) {\x};
  }
\end{tikzpicture}
   (width ~106 mm)
\begin{tikzpicture}[x=10mm, y=1mm, font=\footnotesize\sffamily,
  t1/.style={draw, fill=ln-accent!25},
  t2/.style={draw, fill=ln-box},
  t3/.style={draw, fill=white},
  lab/.style={anchor=east, align=right}]
  % FCFS
  \node[lab] at (-0.2,25) {FCFS};
  \draw[t1] (0,22) rectangle (5,28); \node at (2.5,25) {T1};
  \draw[t2] (5,22) rectangle (7,28); \node at (6,25) {T2};
  \draw[t3] (7,22) rectangle (8,28); \node at (7.5,25) {T3};
  % SJF
  \node[lab] at (-0.2,14) {SJF};
  \draw[t3] (0,11) rectangle (1,17); \node at (0.5,14) {T3};
  \draw[t2] (1,11) rectangle (3,17); \node at (2,14) {T2};
  \draw[t1] (3,11) rectangle (8,17); \node at (5.5,14) {T1};
  % round robin, timeslice 1 s
  \node[lab] at (-0.2,3) {\iflangja{RR（1\,s）}{RR (1\,s)}};
  \draw[t1] (0,0) rectangle (1,6); \node at (0.5,3) {T1};
  \draw[t2] (1,0) rectangle (2,6); \node at (1.5,3) {T2};
  \draw[t3] (2,0) rectangle (3,6); \node at (2.5,3) {T3};
  \draw[t1] (3,0) rectangle (4,6); \node at (3.5,3) {T1};
  \draw[t2] (4,0) rectangle (5,6); \node at (4.5,3) {T2};
  \draw[t1] (5,0) rectangle (8,6); \node at (6.5,3) {T1};
  \draw[densely dotted] (6,0) -- (6,6) (7,0) -- (7,6);
  % time axis
  \draw[->] (0,-3) -- (8.6,-3);
  \node[anchor=west] at (8.6,-3) {$t$\,(s)};
  \foreach \x in {0,...,8} {
    \draw (\x,-3) -- (\x,-4.2);
    \node[below, font=\scriptsize\sffamily] at (\x,-4.2) {\x};
  }
\end{tikzpicture}

  \caption{時刻 0 に T1，T2，T3 の順で到着し，実行時間がそれぞれ
    \SI{5}{\second}，\SI{2}{\second}，\SI{1}{\second} である3つのタスクを，
    FCFS，SJF，およびタイムスライス \SI{1}{\second} のラウンドロビン
    （\cref{sec:l07-rr}）でスケジュールした場合．点線は同じタスクの連続
    するタイムスライスの区切りを示す．}
  \label{fig:l07-gantt}
\end{figure}

\begin{example}
  \label{ex:l07-rtc}
  3つのタスクが時刻 0 に T1，T2，T3 の順で到着し，実行時間はそれぞれ
  \SI{5}{\second}，\SI{2}{\second}，\SI{1}{\second} であるとする
  （\cref{fig:l07-gantt}）．FCFS ではそれらは 5，7，\SI{8}{\second} に
  完了する．スループットは毎秒 $3/8 = 0.375$ タスク，平均完了時間は
  $(5+7+8)/3 \approx \SI{6.67}{\second}$，平均待ち時間は
  $(0+5+7)/3 = \SI{4}{\second}$ である．SJF は T3，T2，T1 の順に実行し，
  それらは 1，3，\SI{8}{\second} に完了する．スループットは毎秒 0.375
  タスクのままだが，平均完了時間は $(1+3+8)/3 = \SI{4}{\second}$ に，
  平均待ち時間は $(0+1+3)/3 \approx \SI{1.33}{\second}$ に下がる．
\end{example}

総仕事量はどちらの順序でも同じであるから，一群全体にわたって測った
スループットは変わらない．しかし短いタスクを先に実行することで，SJF は
一群が終わる前のどの時点でもより多くのタスクを完了させており，2つの平均
をともに下げる．実際，同時に到着し，実行完了型でスケジュールされるタスクに
ついては，これより平均待ち時間の小さい順序は存在しない．

\needspace{7\baselineskip}
\section{プリエンプティブスケジューリング}
\label{sec:l07-preemptive}

実際にはタスクは異なる時刻に到着し，到着したタスクの方が実行中のタスク
より CPU を得るにふさわしいことがある．\source{P3L1，プリエンプティブ
スケジューリング：SJF とプリエンプション，優先度；OSTEP ch.~7--8；
Silberschatz ら ch.~5．}より短いかもしれず，より重要かもしれない．実行
完了型のスケジューラは，実行中のタスクが完了するまで，それがどれほど長く
かかろうとも待たせる．

\begin{definition}[プリエンプティブスケジューリング]
  \term[ぷりえんぷてぃぶすけじゅーりんぐ]{プリエンプティブスケジュー%
  リング}[preemptive scheduling]では，スケジューラは実行中のタスクを完了
  前に中断し，CPU を別のタスクに与えることができる．
\end{definition}

SJF は次のようにしてプリエンプティブになる．タスクが到着するたびに，
スケジューラはその実行時間と，実行中のタスクの残りの実行時間とを比較する．
新しいタスクの方が短ければ，実行中のタスクはプリエンプトされてランキュー
に戻され，新しいタスクが実行される．\margin{この方針は\emph{最短残余時間
優先}とも呼ばれる．OSTEP はこれを\emph{最短完了時間優先}（STCF）と呼ぶ．}

タスクの実行時間があらかじめ分かっていることはまれである．そこで
スケジューラは，そのタスクの履歴から，最後に CPU 上にあったときの実行
時間や，直近 $n$ 回の実行時間の窓平均といった経験則を用いて，それを推定
する．\margin{Silberschatz らは指数平均
$\tau_{n+1} = \alpha t_n + (1-\alpha)\tau_n$ を用いる．$\alpha = 1/2$
とすれば保持すべきは $\tau_n$ だけであり，古い実行区間ほど重みが半分ずつ
になる．}

\subsection{優先度スケジューリング}

タスクは重要さの点でも異なる．オペレーティングシステムにとって不可欠な
仕事を行うカーネルレベルスレッドは，ユーザプロセスのスレッドより重要で
ありうるし，ユーザプロセスの中にも他より緊急なものがある．

\begin{definition}[優先度]
  タスクの\term[ゆうせんど]{優先度}[priority]とは，スケジューラにとっての
  そのタスクの重要さを表す値である．
  \term[ゆうせんどすけじゅーりんぐ]{優先度スケジューリング}%
  [priority scheduling]を用いるスケジューラは，実行可能なタスクのうち
  優先度が最も高いものを次に実行し，より高い優先度のタスクが実行可能に
  なったときには実行中のタスクをプリエンプトする．
\end{definition}

優先度スケジューラのランキューは，優先度の水準ごとに1つのキューを持ち，
空でない最上位の水準の先頭のタスクを選ぶ形にできる．あるいは優先度で
順序付けた木にすることもできる．

\begin{example}
  \label{ex:l07-preemptive}
  3つのタスクが次のように到着するとする（\cref{fig:l07-preemptive}）．

  \begin{center}\small
  \begin{tabular}{@{}lccc@{}}
    \toprule
    タスク & 到着（s） & 実行時間（s） & 優先度 \\
    \midrule
    T1 & 0 & 7 & 中 \\
    T2 & 1 & 2 & 低 \\
    T3 & 2 & 3 & 高 \\
    \bottomrule
  \end{tabular}
  \end{center}

  プリエンプティブ SJF では，T1 は \SI{1}{\second} に T2 が到着するまで
  実行される．T2 は T1 に残る \SI{6}{\second} より短いのでこれを
  プリエンプトする．\SI{2}{\second} に T3 が到着したとき，T2 の残りは
  \SI{1}{\second} しかないので T2 が実行を続け，\SI{3}{\second} に完了
  する．T3 は T1 の残りより短いので次に実行され，\SI{6}{\second} に完了
  する．T1 は \SI{12}{\second} に完了する．到着から測ると各タスクは 12，
  2，\SI{4}{\second} を要し，平均は \SI{6}{\second} である．プリエンプ
  ティブな優先度スケジューリングでは，T2 は T1 をプリエンプトしないが
  T3 はプリエンプトする．T3 が 2 から \SI{5}{\second} まで実行され，
  次に T1 が \SI{10}{\second} に完了し，T2 は \SI{12}{\second} になって
  ようやく完了する．
\end{example}

\begin{figure}[htbp]
  \centering
  % Preemptive SJF and preemptive priority scheduling of three tasks:
% T1 arrives at 0 (7 s, medium), T2 at 1 (2 s, low), T3 at 2 (3 s, high).
% Usage: % Preemptive SJF and preemptive priority scheduling of three tasks:
% T1 arrives at 0 (7 s, medium), T2 at 1 (2 s, low), T3 at 2 (3 s, high).
% Usage: % Preemptive SJF and preemptive priority scheduling of three tasks:
% T1 arrives at 0 (7 s, medium), T2 at 1 (2 s, low), T3 at 2 (3 s, high).
% Usage: \input{figures/l07-preemptive}   (width ~100 mm)
\begin{tikzpicture}[x=7.5mm, y=1mm, font=\footnotesize\sffamily,
  >={Stealth[length=1.5mm]},
  t1/.style={draw, fill=ln-accent!25},
  t2/.style={draw, fill=ln-box},
  t3/.style={draw, fill=white},
  ttl/.style={anchor=south west, font=\footnotesize\sffamily\bfseries, inner sep=1pt}]
  % preemptive SJF
  \node[ttl] at (0,26.5) {\iflangja{(a) プリエンプティブ SJF}{(a) preemptive SJF}};
  \draw[t1] (0,20) rectangle (1,26);  \node at (0.5,23) {T1};
  \draw[t2] (1,20) rectangle (3,26);  \node at (2,23) {T2};
  \draw[t3] (3,20) rectangle (6,26);  \node at (4.5,23) {T3};
  \draw[t1] (6,20) rectangle (12,26); \node at (9,23) {T1};
  % priority
  \node[ttl] at (0,10.5) {\iflangja{(b) プリエンプティブ優先度スケジューリング}{(b) preemptive priority scheduling}};
  \draw[t1] (0,4) rectangle (2,10);   \node at (1,7) {T1};
  \draw[t3] (2,4) rectangle (5,10);   \node at (3.5,7) {T3};
  \draw[t1] (5,4) rectangle (10,10);  \node at (7.5,7) {T1};
  \draw[t2] (10,4) rectangle (12,10); \node at (11,7) {T2};
  % time axis
  \draw[->] (0,0) -- (12.8,0);
  \node[anchor=west] at (12.8,0) {$t$\,(s)};
  \foreach \x in {0,...,12} {
    \draw (\x,0) -- (\x,-1.2);
    \node[below, font=\scriptsize\sffamily] at (\x,-1.2) {\x};
  }
  % arrivals
  \node[anchor=east, font=\scriptsize\sffamily] at (-0.2,-9) {\iflangja{到着}{arrival}};
  \foreach \x/\t in {0/T1, 1/T2, 2/T3} {
    \draw[->] (\x,-10.5) -- (\x,-5.5);
    \node[below, font=\scriptsize\sffamily, inner sep=1pt] at (\x,-10.5) {\t};
  }
\end{tikzpicture}
   (width ~100 mm)
\begin{tikzpicture}[x=7.5mm, y=1mm, font=\footnotesize\sffamily,
  >={Stealth[length=1.5mm]},
  t1/.style={draw, fill=ln-accent!25},
  t2/.style={draw, fill=ln-box},
  t3/.style={draw, fill=white},
  ttl/.style={anchor=south west, font=\footnotesize\sffamily\bfseries, inner sep=1pt}]
  % preemptive SJF
  \node[ttl] at (0,26.5) {\iflangja{(a) プリエンプティブ SJF}{(a) preemptive SJF}};
  \draw[t1] (0,20) rectangle (1,26);  \node at (0.5,23) {T1};
  \draw[t2] (1,20) rectangle (3,26);  \node at (2,23) {T2};
  \draw[t3] (3,20) rectangle (6,26);  \node at (4.5,23) {T3};
  \draw[t1] (6,20) rectangle (12,26); \node at (9,23) {T1};
  % priority
  \node[ttl] at (0,10.5) {\iflangja{(b) プリエンプティブ優先度スケジューリング}{(b) preemptive priority scheduling}};
  \draw[t1] (0,4) rectangle (2,10);   \node at (1,7) {T1};
  \draw[t3] (2,4) rectangle (5,10);   \node at (3.5,7) {T3};
  \draw[t1] (5,4) rectangle (10,10);  \node at (7.5,7) {T1};
  \draw[t2] (10,4) rectangle (12,10); \node at (11,7) {T2};
  % time axis
  \draw[->] (0,0) -- (12.8,0);
  \node[anchor=west] at (12.8,0) {$t$\,(s)};
  \foreach \x in {0,...,12} {
    \draw (\x,0) -- (\x,-1.2);
    \node[below, font=\scriptsize\sffamily] at (\x,-1.2) {\x};
  }
  % arrivals
  \node[anchor=east, font=\scriptsize\sffamily] at (-0.2,-9) {\iflangja{到着}{arrival}};
  \foreach \x/\t in {0/T1, 1/T2, 2/T3} {
    \draw[->] (\x,-10.5) -- (\x,-5.5);
    \node[below, font=\scriptsize\sffamily, inner sep=1pt] at (\x,-10.5) {\t};
  }
\end{tikzpicture}
   (width ~100 mm)
\begin{tikzpicture}[x=7.5mm, y=1mm, font=\footnotesize\sffamily,
  >={Stealth[length=1.5mm]},
  t1/.style={draw, fill=ln-accent!25},
  t2/.style={draw, fill=ln-box},
  t3/.style={draw, fill=white},
  ttl/.style={anchor=south west, font=\footnotesize\sffamily\bfseries, inner sep=1pt}]
  % preemptive SJF
  \node[ttl] at (0,26.5) {\iflangja{(a) プリエンプティブ SJF}{(a) preemptive SJF}};
  \draw[t1] (0,20) rectangle (1,26);  \node at (0.5,23) {T1};
  \draw[t2] (1,20) rectangle (3,26);  \node at (2,23) {T2};
  \draw[t3] (3,20) rectangle (6,26);  \node at (4.5,23) {T3};
  \draw[t1] (6,20) rectangle (12,26); \node at (9,23) {T1};
  % priority
  \node[ttl] at (0,10.5) {\iflangja{(b) プリエンプティブ優先度スケジューリング}{(b) preemptive priority scheduling}};
  \draw[t1] (0,4) rectangle (2,10);   \node at (1,7) {T1};
  \draw[t3] (2,4) rectangle (5,10);   \node at (3.5,7) {T3};
  \draw[t1] (5,4) rectangle (10,10);  \node at (7.5,7) {T1};
  \draw[t2] (10,4) rectangle (12,10); \node at (11,7) {T2};
  % time axis
  \draw[->] (0,0) -- (12.8,0);
  \node[anchor=west] at (12.8,0) {$t$\,(s)};
  \foreach \x in {0,...,12} {
    \draw (\x,0) -- (\x,-1.2);
    \node[below, font=\scriptsize\sffamily] at (\x,-1.2) {\x};
  }
  % arrivals
  \node[anchor=east, font=\scriptsize\sffamily] at (-0.2,-9) {\iflangja{到着}{arrival}};
  \foreach \x/\t in {0/T1, 1/T2, 2/T3} {
    \draw[->] (\x,-10.5) -- (\x,-5.5);
    \node[below, font=\scriptsize\sffamily, inner sep=1pt] at (\x,-10.5) {\t};
  }
\end{tikzpicture}

  \caption{\cref{ex:l07-preemptive}のタスクを (a) プリエンプティブ SJF
    と (b) プリエンプティブな優先度スケジューリングで実行した場合．}
  \label{fig:l07-preemptive}
\end{figure}

\subsection{スタベーションとエージング}

優先度スケジューリングには危険がある．より高い優先度のタスクが到着し続け
る限り，低い優先度のタスクはランキューに留まり続ける．

\begin{definition}[スタベーション]
  スケジューラが常に他のタスクを優先するために，実行可能でありながら CPU
  を無期限に与えられないタスクは，\term[すたべーしょん]{スタベーション}%
  [starvation]に陥る．
\end{definition}

その対策が\term[えーじんぐ]{エージング}[priority aging]である．
スケジューラが用いる優先度は，タスク本来の優先度だけではなく，本来の
優先度と，そのタスクがランキューで過ごした時間との関数である．タスクが
待つ時間が長いほどその実効的な優先度は高くなり，やがて到着するタスクの
優先度を上回ってそのタスクが実行される．こうしてエージングはスタベー
ションを防ぐ．

\needspace{7\baselineskip}
\section{優先度逆転}
\label{sec:l07-inversion}

優先度は同期と相互に作用する．\source{P3L1，優先度逆転；Silberschatz ら
ch.~5--6．}優先度が $\text{T1} > \text{T2} > \text{T3}$ である3つの
タスクを考える．このうち T1 と T3 は同じミューテックス \texttt{m} を
ロックする（\cref{fig:l07-inversion}a）．T3 が最初に到着し，時刻 1 に
\texttt{m} をロックする．時刻 2 に T2 が到着して T3 をプリエンプトし，
時刻 3 に T1 が到着して T2 をプリエンプトする．時刻 4 に T1 は
\texttt{m} をロックしようとするが，これは T3 が保持しているのでブロック
する．残る2つのタスクのうち優先度が高いのは T2 であるから，T3 が実行を
再開できるようになる前に T2 が最後まで実行される．その後ようやく
T3 はクリティカルセクションを終えて \texttt{m} をアンロックし，T1 が先へ
進める．T2 は T1 より早く時刻 6 に完了し，T1 は時刻 9 に完了する．

\begin{definition}[優先度逆転]
  \term[ゆうせんどぎゃくてん]{優先度逆転}[priority inversion]とは，高い
  優先度のタスクが，より低い優先度のタスクが保持する資源を待っている
  一方で，その資源を必要としない中間の優先度のタスクがその間に実行される
  ことである．高い優先度のタスクは，あたかも優先度が逆転しているかの
  ように遅れる．
\end{definition}

\begin{figure}[htbp]
  \centering
  % Priority inversion and priority inheritance.  T1 (high) arrives at 3,
% T2 (medium) at 2, T3 (low) at 0; T1 and T3 share mutex m.
% Usage: % Priority inversion and priority inheritance.  T1 (high) arrives at 3,
% T2 (medium) at 2, T3 (low) at 0; T1 and T3 share mutex m.
% Usage: % Priority inversion and priority inheritance.  T1 (high) arrives at 3,
% T2 (medium) at 2, T3 (low) at 0; T1 and T3 share mutex m.
% Usage: \input{figures/l07-inversion}   (width ~112 mm)
\begin{tikzpicture}[x=9mm, y=1mm, font=\footnotesize\sffamily,
  run/.style={draw, fill=ln-box},
  cs/.style={draw, fill=ln-accent!40},
  boost/.style={draw=ln-caution, line width=1.2pt, fill=ln-accent!40},
  blk/.style={densely dashed, thick, ln-caution},
  lab/.style={anchor=east},
  ttl/.style={anchor=south west, font=\footnotesize\sffamily\bfseries, inner sep=1pt}]
  % ---------- (a) without priority inheritance
  \node[ttl] at (0,54) {\iflangja{(a) 優先度継承なし}{(a) without priority inheritance}};
  \node[lab] at (-0.1,50.5) {\iflangja{T1（高）}{T1 (high)}};
  \node[lab] at (-0.1,42.5) {\iflangja{T2（中）}{T2 (medium)}};
  \node[lab] at (-0.1,34.5) {\iflangja{T3（低）}{T3 (low)}};
  \draw[run] (3,48) rectangle (4,53);
  \draw[blk] (4,50.5) -- (7,50.5);
  \draw[cs]  (7,48) rectangle (8,53);
  \draw[run] (8,48) rectangle (9,53);
  \draw[run] (2,40) rectangle (3,45);
  \draw[run] (4,40) rectangle (6,45);
  \draw[run] (0,32) rectangle (1,37);
  \draw[cs]  (1,32) rectangle (2,37);
  \draw[cs]  (6,32) rectangle (7,37);
  \draw[run] (9,32) rectangle (10,37);
  \draw (0,29) -- (10.3,29);
  \foreach \x in {0,...,10} {
    \draw (\x,29) -- (\x,28);
    \node[below, font=\scriptsize\sffamily, inner sep=1pt] at (\x,28) {\x};
  }
  % ---------- (b) with priority inheritance
  \node[ttl] at (0,19) {\iflangja{(b) 優先度継承あり}{(b) with priority inheritance}};
  \node[lab] at (-0.1,15.5) {\iflangja{T1（高）}{T1 (high)}};
  \node[lab] at (-0.1,7.5)  {\iflangja{T2（中）}{T2 (medium)}};
  \node[lab] at (-0.1,-0.5) {\iflangja{T3（低）}{T3 (low)}};
  \draw[run] (3,13) rectangle (4,18);
  \draw[blk] (4,15.5) -- (5,15.5);
  \draw[cs]  (5,13) rectangle (6,18);
  \draw[run] (6,13) rectangle (7,18);
  \draw[run] (2,5) rectangle (3,10);
  \draw[run] (7,5) rectangle (9,10);
  \draw[run] (0,-3) rectangle (1,2);
  \draw[cs]  (1,-3) rectangle (2,2);
  \draw[boost] (4,-3) rectangle (5,2);
  \node[anchor=west, font=\scriptsize\sffamily, text=ln-caution] at (5.1,-0.5)
    {\iflangja{優先度を「高」に引き上げ}{priority raised to high}};
  \draw[run] (9,-3) rectangle (10,2);
  \draw (0,-6) -- (10.3,-6);
  \foreach \x in {0,...,10} {
    \draw (\x,-6) -- (\x,-7);
    \node[below, font=\scriptsize\sffamily, inner sep=1pt] at (\x,-7) {\x};
  }
  \node[anchor=west, font=\scriptsize\sffamily] at (10.3,-6) {$t$};
  % ---------- legend
  \begin{scope}[yshift=-17mm, font=\scriptsize\sffamily]
    \draw[run] (0,-1.5) rectangle (0.5,1.5);
    \node[anchor=west] at (0.55,0) {\iflangja{実行}{runs}};
    \draw[cs] (2,-1.5) rectangle (2.5,1.5);
    \node[anchor=west] at (2.55,0) {\iflangja{m を保持して実行}{runs holding m}};
    % The Japanese labels are wider: shift the last two entries so that the
    % "inherited priority" swatch does not cover the end of "blocked on m".
    \draw[blk] (\iflangja{5.0}{5.3},0) -- (\iflangja{5.6}{6},0);
    \node[anchor=west] at (\iflangja{5.65}{6.05},0) {\iflangja{m 待ちでブロック}{blocked on m}};
    \draw[boost] (\iflangja{8.45}{8.2},-1.5) rectangle (\iflangja{8.95}{8.7},1.5);
    \node[anchor=west] at (\iflangja{9.0}{8.75},0) {\iflangja{継承した優先度}{inherited priority}};
  \end{scope}
\end{tikzpicture}
   (width ~112 mm)
\begin{tikzpicture}[x=9mm, y=1mm, font=\footnotesize\sffamily,
  run/.style={draw, fill=ln-box},
  cs/.style={draw, fill=ln-accent!40},
  boost/.style={draw=ln-caution, line width=1.2pt, fill=ln-accent!40},
  blk/.style={densely dashed, thick, ln-caution},
  lab/.style={anchor=east},
  ttl/.style={anchor=south west, font=\footnotesize\sffamily\bfseries, inner sep=1pt}]
  % ---------- (a) without priority inheritance
  \node[ttl] at (0,54) {\iflangja{(a) 優先度継承なし}{(a) without priority inheritance}};
  \node[lab] at (-0.1,50.5) {\iflangja{T1（高）}{T1 (high)}};
  \node[lab] at (-0.1,42.5) {\iflangja{T2（中）}{T2 (medium)}};
  \node[lab] at (-0.1,34.5) {\iflangja{T3（低）}{T3 (low)}};
  \draw[run] (3,48) rectangle (4,53);
  \draw[blk] (4,50.5) -- (7,50.5);
  \draw[cs]  (7,48) rectangle (8,53);
  \draw[run] (8,48) rectangle (9,53);
  \draw[run] (2,40) rectangle (3,45);
  \draw[run] (4,40) rectangle (6,45);
  \draw[run] (0,32) rectangle (1,37);
  \draw[cs]  (1,32) rectangle (2,37);
  \draw[cs]  (6,32) rectangle (7,37);
  \draw[run] (9,32) rectangle (10,37);
  \draw (0,29) -- (10.3,29);
  \foreach \x in {0,...,10} {
    \draw (\x,29) -- (\x,28);
    \node[below, font=\scriptsize\sffamily, inner sep=1pt] at (\x,28) {\x};
  }
  % ---------- (b) with priority inheritance
  \node[ttl] at (0,19) {\iflangja{(b) 優先度継承あり}{(b) with priority inheritance}};
  \node[lab] at (-0.1,15.5) {\iflangja{T1（高）}{T1 (high)}};
  \node[lab] at (-0.1,7.5)  {\iflangja{T2（中）}{T2 (medium)}};
  \node[lab] at (-0.1,-0.5) {\iflangja{T3（低）}{T3 (low)}};
  \draw[run] (3,13) rectangle (4,18);
  \draw[blk] (4,15.5) -- (5,15.5);
  \draw[cs]  (5,13) rectangle (6,18);
  \draw[run] (6,13) rectangle (7,18);
  \draw[run] (2,5) rectangle (3,10);
  \draw[run] (7,5) rectangle (9,10);
  \draw[run] (0,-3) rectangle (1,2);
  \draw[cs]  (1,-3) rectangle (2,2);
  \draw[boost] (4,-3) rectangle (5,2);
  \node[anchor=west, font=\scriptsize\sffamily, text=ln-caution] at (5.1,-0.5)
    {\iflangja{優先度を「高」に引き上げ}{priority raised to high}};
  \draw[run] (9,-3) rectangle (10,2);
  \draw (0,-6) -- (10.3,-6);
  \foreach \x in {0,...,10} {
    \draw (\x,-6) -- (\x,-7);
    \node[below, font=\scriptsize\sffamily, inner sep=1pt] at (\x,-7) {\x};
  }
  \node[anchor=west, font=\scriptsize\sffamily] at (10.3,-6) {$t$};
  % ---------- legend
  \begin{scope}[yshift=-17mm, font=\scriptsize\sffamily]
    \draw[run] (0,-1.5) rectangle (0.5,1.5);
    \node[anchor=west] at (0.55,0) {\iflangja{実行}{runs}};
    \draw[cs] (2,-1.5) rectangle (2.5,1.5);
    \node[anchor=west] at (2.55,0) {\iflangja{m を保持して実行}{runs holding m}};
    % The Japanese labels are wider: shift the last two entries so that the
    % "inherited priority" swatch does not cover the end of "blocked on m".
    \draw[blk] (\iflangja{5.0}{5.3},0) -- (\iflangja{5.6}{6},0);
    \node[anchor=west] at (\iflangja{5.65}{6.05},0) {\iflangja{m 待ちでブロック}{blocked on m}};
    \draw[boost] (\iflangja{8.45}{8.2},-1.5) rectangle (\iflangja{8.95}{8.7},1.5);
    \node[anchor=west] at (\iflangja{9.0}{8.75},0) {\iflangja{継承した優先度}{inherited priority}};
  \end{scope}
\end{tikzpicture}
   (width ~112 mm)
\begin{tikzpicture}[x=9mm, y=1mm, font=\footnotesize\sffamily,
  run/.style={draw, fill=ln-box},
  cs/.style={draw, fill=ln-accent!40},
  boost/.style={draw=ln-caution, line width=1.2pt, fill=ln-accent!40},
  blk/.style={densely dashed, thick, ln-caution},
  lab/.style={anchor=east},
  ttl/.style={anchor=south west, font=\footnotesize\sffamily\bfseries, inner sep=1pt}]
  % ---------- (a) without priority inheritance
  \node[ttl] at (0,54) {\iflangja{(a) 優先度継承なし}{(a) without priority inheritance}};
  \node[lab] at (-0.1,50.5) {\iflangja{T1（高）}{T1 (high)}};
  \node[lab] at (-0.1,42.5) {\iflangja{T2（中）}{T2 (medium)}};
  \node[lab] at (-0.1,34.5) {\iflangja{T3（低）}{T3 (low)}};
  \draw[run] (3,48) rectangle (4,53);
  \draw[blk] (4,50.5) -- (7,50.5);
  \draw[cs]  (7,48) rectangle (8,53);
  \draw[run] (8,48) rectangle (9,53);
  \draw[run] (2,40) rectangle (3,45);
  \draw[run] (4,40) rectangle (6,45);
  \draw[run] (0,32) rectangle (1,37);
  \draw[cs]  (1,32) rectangle (2,37);
  \draw[cs]  (6,32) rectangle (7,37);
  \draw[run] (9,32) rectangle (10,37);
  \draw (0,29) -- (10.3,29);
  \foreach \x in {0,...,10} {
    \draw (\x,29) -- (\x,28);
    \node[below, font=\scriptsize\sffamily, inner sep=1pt] at (\x,28) {\x};
  }
  % ---------- (b) with priority inheritance
  \node[ttl] at (0,19) {\iflangja{(b) 優先度継承あり}{(b) with priority inheritance}};
  \node[lab] at (-0.1,15.5) {\iflangja{T1（高）}{T1 (high)}};
  \node[lab] at (-0.1,7.5)  {\iflangja{T2（中）}{T2 (medium)}};
  \node[lab] at (-0.1,-0.5) {\iflangja{T3（低）}{T3 (low)}};
  \draw[run] (3,13) rectangle (4,18);
  \draw[blk] (4,15.5) -- (5,15.5);
  \draw[cs]  (5,13) rectangle (6,18);
  \draw[run] (6,13) rectangle (7,18);
  \draw[run] (2,5) rectangle (3,10);
  \draw[run] (7,5) rectangle (9,10);
  \draw[run] (0,-3) rectangle (1,2);
  \draw[cs]  (1,-3) rectangle (2,2);
  \draw[boost] (4,-3) rectangle (5,2);
  \node[anchor=west, font=\scriptsize\sffamily, text=ln-caution] at (5.1,-0.5)
    {\iflangja{優先度を「高」に引き上げ}{priority raised to high}};
  \draw[run] (9,-3) rectangle (10,2);
  \draw (0,-6) -- (10.3,-6);
  \foreach \x in {0,...,10} {
    \draw (\x,-6) -- (\x,-7);
    \node[below, font=\scriptsize\sffamily, inner sep=1pt] at (\x,-7) {\x};
  }
  \node[anchor=west, font=\scriptsize\sffamily] at (10.3,-6) {$t$};
  % ---------- legend
  \begin{scope}[yshift=-17mm, font=\scriptsize\sffamily]
    \draw[run] (0,-1.5) rectangle (0.5,1.5);
    \node[anchor=west] at (0.55,0) {\iflangja{実行}{runs}};
    \draw[cs] (2,-1.5) rectangle (2.5,1.5);
    \node[anchor=west] at (2.55,0) {\iflangja{m を保持して実行}{runs holding m}};
    % The Japanese labels are wider: shift the last two entries so that the
    % "inherited priority" swatch does not cover the end of "blocked on m".
    \draw[blk] (\iflangja{5.0}{5.3},0) -- (\iflangja{5.6}{6},0);
    \node[anchor=west] at (\iflangja{5.65}{6.05},0) {\iflangja{m 待ちでブロック}{blocked on m}};
    \draw[boost] (\iflangja{8.45}{8.2},-1.5) rectangle (\iflangja{8.95}{8.7},1.5);
    \node[anchor=west] at (\iflangja{9.0}{8.75},0) {\iflangja{継承した優先度}{inherited priority}};
  \end{scope}
\end{tikzpicture}

  \caption{T1 と T3 はミューテックス \texttt{m} を共有する．(a) T1 が
    \texttt{m} でブロックしている間に，中優先度の T2 が T1 に先んじて
    完了まで実行される．(b) 優先度継承があると，T3 は \texttt{m} を
    アンロックするまで T1 の優先度で実行され，T1 は T2 より先に完了
    する．}
  \label{fig:l07-inversion}
\end{figure}

解決策は，ミューテックスの所有者の優先度を一時的に引き上げることである
（\cref{fig:l07-inversion}b）．T1 が \texttt{m} でブロックすると，T3 に
T1 の優先度が与えられ，T2 ではなく T3 が実行されてクリティカルセクション
を終える．T3 が \texttt{m} をアンロックすると，その優先度は本来のものに
戻され，T1 がミューテックスを獲得して実行され，T1 は T2 に先んじて時刻 7
に完了する．

\begin{definition}[優先度継承]
  \term[ゆうせんどけいしょう]{優先度継承}[priority inheritance]では，
  ロックの所有者は，そのロックでブロックしているタスクのうち最も優先度の
  高いものの優先度で一時的に実行され，ロックを解放すると自分本来の優先度
  に戻る．
\end{definition}

\needspace{9\baselineskip}
所有者の優先度を引き上げるには，ロック操作が誰が所有者であるかを知らな
ければならない．これも，ミューテックスのデータ構造に所有者を記録して
おく理由の1つである（第5章）．\sidenote{1997年7月，火星に着陸した
Mars Pathfinder は再起動を繰り返した．優先度の低い気象観測タスクが共有情報
バスのミューテックスを保持し，優先度の高いバスのタスクがそれでブロックし，
代わりに中優先度の通信処理が実行され，最後にウォッチドッグタイマが探査機を
リセットしたのである．JPL は地上で優先度逆転を再現し，そのミューテックスで
優先度継承を有効にする修正を送信した．}%
POSIX スレッド（第4章）は，優先度継承を
ミューテックスのプロトコル属性として提供する．
\begin{lstlisting}[language=C, numbers=none]
pthread_mutexattr_t attr;
pthread_mutex_t m;

pthread_mutexattr_init(&attr);
pthread_mutexattr_setprotocol(&attr, PTHREAD_PRIO_INHERIT);
pthread_mutex_init(&m, &attr);
\end{lstlisting}

\begin{keypoint}
  プリエンプションによって，スケジューラは短いタスクや重要なタスクを先に
  処理できる．優先度は2つの危険をもたらす．1つはスタベーションであり，
  エージングによって回避される．もう1つは優先度逆転であり，ロックの所有者
  にそのロックでブロックしているタスクの優先度を継承させることで回避される．%
  \caution{POSIX は優先度上限プロトコル
  \texttt{PTHREAD\_PRIO\_PROTECT} も定義している．所有者は，ブロックして
  いるタスクの有無によらず，ミューテックスに宣言された上限の優先度で実行
  される．講義は継承のみを扱う．}
\end{keypoint}

\needspace{7\baselineskip}
\section{ラウンドロビンスケジューリングとタイムシェアリング}
\label{sec:l07-rr}

実行時間が分からないタスクのためのスケジューラは，単純にタスクを順番に
実行させればよい．\source{P3L1，ラウンドロビンスケジューリング，
タイムシェアリングとタイムスライス；OSTEP ch.~7．}このようなスケジューラ
には，実行時間の推定も，基本的な形では優先度も必要ない．

\begin{definition}[ラウンドロビン]
  \term[らうんどろびんすけじゅーりんぐ]{ラウンドロビンスケジューリング}%
  [round robin scheduling]（RR）を用いるスケジューラは，FCFS と同様に
  ランキューの先頭のタスクを取り出す．FCFS と異なり，タスクは完了まで
  実行される必要はない．例えば I/O を待つために CPU を手放すことができ，
  再び実行可能になったときにはランキューの末尾に置かれる．
\end{definition}

ラウンドロビンは優先度と組み合わせられる．同じ優先度のタスクが順番に
実行され，より高い優先度のタスクが実行可能になれば，実行中のタスクを
プリエンプトする．最も一般的な形では，ラウンドロビンは各タスクが実行して
よい時間を制限することで，タスクが自ら CPU を手放さない場合にもタスクを
交互に実行する．

第2章のタイムスライス\index{たいむすらいす@タイムスライス}は，タスクに
与えられる，中断されない時間の最大値であり，タイムクォンタムとも呼ばれる．%
\margin{Linux は今もリアルタイムのラウンドロビンクラスには固定のタイム
クォンタムを用いる．既定は \SI{100}{\milli\second} であり，sysctl の
\texttt{sched\_rr\_timeslice\_ms} で設定する．}%
タスクはタイムスライスより短い時間しか実行しないこともある．I/O や同期
機構を待ってブロックし，対応するキューに置かれるかもしれないし，より高い
優先度のタスクが実行可能になってプリエンプトするかもしれない．タイム
スライスが満了すると，実行中のタスクはプリエンプトされてランキューの末尾
に戻る．

\begin{definition}[タイムシェアリング]
  タイムスライスによってタスクを交互に実行することを，CPU の
  \term[たいむしぇありんぐ]{タイムシェアリング}[timesharing]という．
  どのタスクも一度に高々1タイムスライスしか実行されないので，CPU 集約的
  なタスクもタイムスライスが満了すればプリエンプトされ，完了するまで CPU
  を占有し続けることはできない．
\end{definition}

\begin{example}
  \cref{ex:l07-rtc}のタスクを，コンテキストスイッチの費用を無視して，
  タイムスライス \SI{1}{\second} のラウンドロビンで実行する
  （\cref{fig:l07-gantt}）．順序は T1，T2，T3，T1，T2，T1 であり，T3 は
  \SI{3}{\second}，T2 は \SI{5}{\second}，T1 は最後の3タイムスライスを
  単独で実行した後 \SI{8}{\second} に完了する．指標は \cref{tab:l07-rr}
  で比較する．実行時間を一切知らなくても，ラウンドロビンは平均完了時間を
  \SI{6.67}{\second} から \SI{5.33}{\second} へ，FCFS から SJF の
  \SI{4.00}{\second} までのちょうど中間まで下げる．また，各タスクが最初に
  実行されるまでに待つのは，自分より前にあるタスクそれぞれについて高々1
  タイムスライス分であるから，その平均待ち時間は SJF より短くさえある．
\end{example}

\begin{table}[htbp]
  \centering
  \caption{\cref{fig:l07-gantt}の3つのスケジュールの指標．}
  \label{tab:l07-rr}
  \small
  \begin{tabular}{@{}lccc@{}}
    \toprule
    & FCFS & SJF & RR（\SI{1}{\second}） \\
    \midrule
    スループット（タスク/秒） & 0.375 & 0.375 & 0.375 \\
    平均完了時間（s） & 6.67 & 4.00 & 5.33 \\
    平均待ち時間（s） & 4.00 & 1.33 & 1.00 \\
    \bottomrule
  \end{tabular}
\end{table}

タイムシェアリングにはいくつかの利点がある．短いタスクがより早く完了する．
実行可能なタスクはどれも限られた時間内に実行されるので，システムの応答性が
高まる．そして，長い I/O 操作を開始するタスクがそれに早く到達するので，
他のタスクが計算している間にその操作を進められる．一方で，先の例が無視した
費用もある．タイムスライスの終わりごとにタイマが実行中のタスクに割り込み，
スケジューラが動作し，コンテキストスイッチが行われる．このオーバーヘッド
を小さく保つには，スケジューラ単独の時間について \cref{eq:l02-useful} が
示したように，タイムスライスはスケジューラとコンテキストスイッチに要する
時間よりはるかに長くなければならない．

\needspace{7\baselineskip}
\section{タイムスライスの長さの選択}
\label{sec:l07-timeslice}

タイムスライスがどれだけ長くあるべきかは，タスクがそれをどう使うかに
依存する．\source{P3L1，タイムスライスはどれだけ長くあるべきか，CPU
バウンドのタイムスライス長，I/O バウンドのタイムスライス長，タイム
スライス長のまとめ；OSTEP ch.~7．}タイムスライスは，タイムシェアリングの
利点とそのオーバーヘッドとを釣り合わせなければならず，その釣り合いは，
主に計算するタスクと，主に I/O を待つタスクとで異なる．

\subsection{CPU バウンドのタスク}

ほとんどの時間を計算に費やし，めったにブロックしないタスクを
\term[CPUばうんどたすく]{CPUバウンドタスク}[CPU-bound task]という．

\begin{example}
  \label{ex:l07-cpu-bound}
  2つの CPU バウンドのタスク T1 と T2 が，それぞれ \SI{6}{\second} の
  CPU 時間を必要とし，時刻 0 に到着するとする．コンテキストスイッチには
  \SI{0.2}{\second} を要する．\caution{\SI{0.2}{\second} は説明のための
  値であって測定値ではない．第2章によればコンテキストスイッチの直接の費用は
  数マイクロ秒であり，実際の系がタイムスライスごとに失うものはこの例より
  はるかに小さい．}%
  タイムスライスが \SI{1}{\second} のとき，
  2つのタスクは 12 回交互に実行され，タイムスライスの間に 11 回の
  コンテキストスイッチが入る．T2 は \SI{1.2}{\second} に開始する．T1 の
  最後のタイムスライスは，10 回のタイムスライスと 10 回のコンテキスト
  スイッチの後，\SI{12}{\second} に開始するので，T1 は \SI{13}{\second}
  に完了し，T2 はもう1回のコンテキストスイッチの後 \SI{14.2}{\second} に
  完了する．タイムスライスが \SI{3}{\second} のときは4回交互に実行され，
  コンテキストスイッチは3回である．T2 は \SI{3.2}{\second} に開始し，T1
  は \SI{9.4}{\second} に，T2 は \SI{12.6}{\second} に完了する．タイム
  スライスが無限のときは，T1 が \SI{6}{\second} に完了するまで実行され，
  T2 は 6.2 から \SI{12.2}{\second} まで実行される．得られる指標を
  \cref{tab:l07-cpu-bound}に挙げる．
\end{example}

\begin{table}[htbp]
  \centering
  \caption{コンテキストスイッチが \SI{0.2}{\second} のときの，
    \SI{6}{\second} の CPU バウンドのタスク2つ
    （\cref{ex:l07-cpu-bound}）．}
  \label{tab:l07-cpu-bound}
  \small
  \begin{tabular}{@{}lccc@{}}
    \toprule
    タイムスライス & スループット（タスク/秒） & 平均待ち時間（s） &
      平均完了時間（s） \\
    \midrule
    \SI{1}{\second}  & $2/14.2 \approx 0.141$ & $(0+1.2)/2 = 0.6$ & $(13+14.2)/2 = 13.6$ \\
    \SI{3}{\second}  & $2/12.6 \approx 0.159$ & $(0+3.2)/2 = 1.6$ & $(9.4+12.6)/2 = 11.0$ \\
    $\infty$         & $2/12.2 \approx 0.164$ & $(0+6.2)/2 = 3.1$ & $(6+12.2)/2 = 9.1$ \\
    \bottomrule
  \end{tabular}
\end{table}

CPU バウンドのタスクにとっては，タイムスライスが長いほどコンテキスト
スイッチに費やされる時間が減るので，スループットと平均完了時間が改善
する．伸びるのは待ち時間だけであり，バッチ計算のような CPU バウンドの
タスクにとって，それはほとんど問題にならない．その利用者が見るのは計算が
完了する時刻であって，開始する時刻ではないからである．極限としての無限の
タイムスライス，すなわち実行完了型が，それらにとっては最良である．

\subsection{I/O バウンドのタスク}

頻繁に I/O でブロックし，I/O 操作の間にわずかしか計算しないタスクを
\term[IOばうんどたすく]{I/Oバウンドタスク}[I/O-bound task]という．すべて
のタスクが I/O バウンドであれば，タイムスライスの長さは，それらの短い
計算の区間を超えてさえいれば，ほとんど問題にならない．どのタスクもタイム
スライスが満了する前に I/O でブロックし，自ら CPU を離れるからである．
長さが問題になるのは，I/O バウンドのタスクが CPU バウンドのタスクと CPU
を共有するときである．

\begin{figure}[htbp]
  \centering
  % A CPU-bound task T1 and an I/O-bound task T2 (5 ms of computation, then
% an I/O operation of 10 ms) under round robin with a context switch of
% 2 ms and timeslices of 5 ms and 30 ms.
% Usage: % A CPU-bound task T1 and an I/O-bound task T2 (5 ms of computation, then
% an I/O operation of 10 ms) under round robin with a context switch of
% 2 ms and timeslices of 5 ms and 30 ms.
% Usage: % A CPU-bound task T1 and an I/O-bound task T2 (5 ms of computation, then
% an I/O operation of 10 ms) under round robin with a context switch of
% 2 ms and timeslices of 5 ms and 30 ms.
% Usage: \input{figures/l07-io-timeslice}   (width ~104 mm)
\begin{tikzpicture}[x=1.1mm, y=1mm, font=\scriptsize\sffamily,
  >={Stealth[length=1.4mm]},
  t1/.style={draw, fill=ln-accent!25},
  t2/.style={draw, fill=white},
  cx/.style={draw, fill=ln-muted},
  io/.style={draw, fill=ln-box},
  lab/.style={anchor=east, font=\footnotesize\sffamily},
  ttl/.style={anchor=south west, font=\footnotesize\sffamily\bfseries, inner sep=1pt}]
  % ---------- (a) timeslice 5 ms
  \node[ttl] at (0,45) {\iflangja{(a) タイムスライス 5\,ms}{(a) timeslice 5\,ms}};
  \node[lab] at (-1,41) {CPU};
  \node[lab] at (-1,31) {I/O};
  \foreach \a in {0,19,38,57} {
    \draw[t2] (\a,38) rectangle (\a+5,44);
    \node at (\a+2.5,41) {T2};
    \draw[cx] (\a+5,38) rectangle (\a+7,44);
    \draw[io] (\a+5,28) rectangle (\a+15,34);
    \node at (\a+10,31) {T2};
  }
  \foreach \a in {7,26,45,64} {
    \draw[t1] (\a,38) rectangle (\a+10,44);
    \draw[densely dotted] (\a+5,38) -- (\a+5,44);
    \node at (\a+2.5,41) {T1};
    \node at (\a+7.5,41) {T1};
    \draw[cx] (\a+10,38) rectangle (\a+12,44);
  }
  \draw[t2] (76,38) -- (80,38) (76,44) -- (80,44) (76,38) -- (76,44);
  \node at (78.2,41) {T2};
  \draw (0,25) -- (82,25);
  \foreach \x in {0,10,...,80} {
    \draw (\x,25) -- (\x,24);
    \node[below, inner sep=1pt] at (\x,24) {\x};
  }
  % ---------- (b) timeslice 30 ms
  \node[ttl] at (0,13) {\iflangja{(b) タイムスライス 30\,ms}{(b) timeslice 30\,ms}};
  \node[lab] at (-1,9) {CPU};
  \node[lab] at (-1,-7) {I/O};
  \foreach \a in {0,39} {
    \draw[t2] (\a,6) rectangle (\a+5,12);
    \node at (\a+2.5,9) {T2};
    \draw[cx] (\a+5,6) rectangle (\a+7,12);
    \draw[io] (\a+5,-10) rectangle (\a+15,-4);
    \node at (\a+10,-7) {T2};
    \draw[t1] (\a+7,6) rectangle (\a+37,12);
    \node at (\a+22,9) {T1};
    \draw[cx] (\a+37,6) rectangle (\a+39,12);
  }
  \draw[t2] (78,6) -- (80,6) (78,12) -- (80,12) (78,6) -- (78,12);
  % T2 is ready but waits for the timeslice of T1 to end
  \draw[<->, ln-caution] (15,2.5) -- (39,2.5);
  \draw[ln-caution, densely dotted] (15,-4) -- (15,2.5) (39,2.5) -- (39,6);
  \node[anchor=north west, text=ln-caution, inner sep=1pt] at (16,1.8)
    {\iflangja{T2 は実行可能だが待つ（24\,ms）}{T2 ready, waiting (24\,ms)}};
  \draw (0,-13) -- (82,-13);
  \foreach \x in {0,10,...,80} {
    \draw (\x,-13) -- (\x,-14);
    \node[below, inner sep=1pt] at (\x,-14) {\x};
  }
  \node[anchor=west] at (82,-13) {ms};
  % ---------- legend
  \begin{scope}[yshift=-24mm]
    \draw[t1] (0,-1.5) rectangle (4,1.5);
    \node[anchor=west] at (4.5,0) {\iflangja{CPU バウンドの T1}{CPU-bound T1}};
    \draw[t2] (28,-1.5) rectangle (32,1.5);
    \node[anchor=west] at (32.5,0) {\iflangja{I/O バウンドの T2}{I/O-bound T2}};
    \draw[cx] (55,-1.5) rectangle (57,1.5);
    \node[anchor=west] at (57.5,0) {\iflangja{コンテキストスイッチ}{context switch}};
  \end{scope}
\end{tikzpicture}
   (width ~104 mm)
\begin{tikzpicture}[x=1.1mm, y=1mm, font=\scriptsize\sffamily,
  >={Stealth[length=1.4mm]},
  t1/.style={draw, fill=ln-accent!25},
  t2/.style={draw, fill=white},
  cx/.style={draw, fill=ln-muted},
  io/.style={draw, fill=ln-box},
  lab/.style={anchor=east, font=\footnotesize\sffamily},
  ttl/.style={anchor=south west, font=\footnotesize\sffamily\bfseries, inner sep=1pt}]
  % ---------- (a) timeslice 5 ms
  \node[ttl] at (0,45) {\iflangja{(a) タイムスライス 5\,ms}{(a) timeslice 5\,ms}};
  \node[lab] at (-1,41) {CPU};
  \node[lab] at (-1,31) {I/O};
  \foreach \a in {0,19,38,57} {
    \draw[t2] (\a,38) rectangle (\a+5,44);
    \node at (\a+2.5,41) {T2};
    \draw[cx] (\a+5,38) rectangle (\a+7,44);
    \draw[io] (\a+5,28) rectangle (\a+15,34);
    \node at (\a+10,31) {T2};
  }
  \foreach \a in {7,26,45,64} {
    \draw[t1] (\a,38) rectangle (\a+10,44);
    \draw[densely dotted] (\a+5,38) -- (\a+5,44);
    \node at (\a+2.5,41) {T1};
    \node at (\a+7.5,41) {T1};
    \draw[cx] (\a+10,38) rectangle (\a+12,44);
  }
  \draw[t2] (76,38) -- (80,38) (76,44) -- (80,44) (76,38) -- (76,44);
  \node at (78.2,41) {T2};
  \draw (0,25) -- (82,25);
  \foreach \x in {0,10,...,80} {
    \draw (\x,25) -- (\x,24);
    \node[below, inner sep=1pt] at (\x,24) {\x};
  }
  % ---------- (b) timeslice 30 ms
  \node[ttl] at (0,13) {\iflangja{(b) タイムスライス 30\,ms}{(b) timeslice 30\,ms}};
  \node[lab] at (-1,9) {CPU};
  \node[lab] at (-1,-7) {I/O};
  \foreach \a in {0,39} {
    \draw[t2] (\a,6) rectangle (\a+5,12);
    \node at (\a+2.5,9) {T2};
    \draw[cx] (\a+5,6) rectangle (\a+7,12);
    \draw[io] (\a+5,-10) rectangle (\a+15,-4);
    \node at (\a+10,-7) {T2};
    \draw[t1] (\a+7,6) rectangle (\a+37,12);
    \node at (\a+22,9) {T1};
    \draw[cx] (\a+37,6) rectangle (\a+39,12);
  }
  \draw[t2] (78,6) -- (80,6) (78,12) -- (80,12) (78,6) -- (78,12);
  % T2 is ready but waits for the timeslice of T1 to end
  \draw[<->, ln-caution] (15,2.5) -- (39,2.5);
  \draw[ln-caution, densely dotted] (15,-4) -- (15,2.5) (39,2.5) -- (39,6);
  \node[anchor=north west, text=ln-caution, inner sep=1pt] at (16,1.8)
    {\iflangja{T2 は実行可能だが待つ（24\,ms）}{T2 ready, waiting (24\,ms)}};
  \draw (0,-13) -- (82,-13);
  \foreach \x in {0,10,...,80} {
    \draw (\x,-13) -- (\x,-14);
    \node[below, inner sep=1pt] at (\x,-14) {\x};
  }
  \node[anchor=west] at (82,-13) {ms};
  % ---------- legend
  \begin{scope}[yshift=-24mm]
    \draw[t1] (0,-1.5) rectangle (4,1.5);
    \node[anchor=west] at (4.5,0) {\iflangja{CPU バウンドの T1}{CPU-bound T1}};
    \draw[t2] (28,-1.5) rectangle (32,1.5);
    \node[anchor=west] at (32.5,0) {\iflangja{I/O バウンドの T2}{I/O-bound T2}};
    \draw[cx] (55,-1.5) rectangle (57,1.5);
    \node[anchor=west] at (57.5,0) {\iflangja{コンテキストスイッチ}{context switch}};
  \end{scope}
\end{tikzpicture}
   (width ~104 mm)
\begin{tikzpicture}[x=1.1mm, y=1mm, font=\scriptsize\sffamily,
  >={Stealth[length=1.4mm]},
  t1/.style={draw, fill=ln-accent!25},
  t2/.style={draw, fill=white},
  cx/.style={draw, fill=ln-muted},
  io/.style={draw, fill=ln-box},
  lab/.style={anchor=east, font=\footnotesize\sffamily},
  ttl/.style={anchor=south west, font=\footnotesize\sffamily\bfseries, inner sep=1pt}]
  % ---------- (a) timeslice 5 ms
  \node[ttl] at (0,45) {\iflangja{(a) タイムスライス 5\,ms}{(a) timeslice 5\,ms}};
  \node[lab] at (-1,41) {CPU};
  \node[lab] at (-1,31) {I/O};
  \foreach \a in {0,19,38,57} {
    \draw[t2] (\a,38) rectangle (\a+5,44);
    \node at (\a+2.5,41) {T2};
    \draw[cx] (\a+5,38) rectangle (\a+7,44);
    \draw[io] (\a+5,28) rectangle (\a+15,34);
    \node at (\a+10,31) {T2};
  }
  \foreach \a in {7,26,45,64} {
    \draw[t1] (\a,38) rectangle (\a+10,44);
    \draw[densely dotted] (\a+5,38) -- (\a+5,44);
    \node at (\a+2.5,41) {T1};
    \node at (\a+7.5,41) {T1};
    \draw[cx] (\a+10,38) rectangle (\a+12,44);
  }
  \draw[t2] (76,38) -- (80,38) (76,44) -- (80,44) (76,38) -- (76,44);
  \node at (78.2,41) {T2};
  \draw (0,25) -- (82,25);
  \foreach \x in {0,10,...,80} {
    \draw (\x,25) -- (\x,24);
    \node[below, inner sep=1pt] at (\x,24) {\x};
  }
  % ---------- (b) timeslice 30 ms
  \node[ttl] at (0,13) {\iflangja{(b) タイムスライス 30\,ms}{(b) timeslice 30\,ms}};
  \node[lab] at (-1,9) {CPU};
  \node[lab] at (-1,-7) {I/O};
  \foreach \a in {0,39} {
    \draw[t2] (\a,6) rectangle (\a+5,12);
    \node at (\a+2.5,9) {T2};
    \draw[cx] (\a+5,6) rectangle (\a+7,12);
    \draw[io] (\a+5,-10) rectangle (\a+15,-4);
    \node at (\a+10,-7) {T2};
    \draw[t1] (\a+7,6) rectangle (\a+37,12);
    \node at (\a+22,9) {T1};
    \draw[cx] (\a+37,6) rectangle (\a+39,12);
  }
  \draw[t2] (78,6) -- (80,6) (78,12) -- (80,12) (78,6) -- (78,12);
  % T2 is ready but waits for the timeslice of T1 to end
  \draw[<->, ln-caution] (15,2.5) -- (39,2.5);
  \draw[ln-caution, densely dotted] (15,-4) -- (15,2.5) (39,2.5) -- (39,6);
  \node[anchor=north west, text=ln-caution, inner sep=1pt] at (16,1.8)
    {\iflangja{T2 は実行可能だが待つ（24\,ms）}{T2 ready, waiting (24\,ms)}};
  \draw (0,-13) -- (82,-13);
  \foreach \x in {0,10,...,80} {
    \draw (\x,-13) -- (\x,-14);
    \node[below, inner sep=1pt] at (\x,-14) {\x};
  }
  \node[anchor=west] at (82,-13) {ms};
  % ---------- legend
  \begin{scope}[yshift=-24mm]
    \draw[t1] (0,-1.5) rectangle (4,1.5);
    \node[anchor=west] at (4.5,0) {\iflangja{CPU バウンドの T1}{CPU-bound T1}};
    \draw[t2] (28,-1.5) rectangle (32,1.5);
    \node[anchor=west] at (32.5,0) {\iflangja{I/O バウンドの T2}{I/O-bound T2}};
    \draw[cx] (55,-1.5) rectangle (57,1.5);
    \node[anchor=west] at (57.5,0) {\iflangja{コンテキストスイッチ}{context switch}};
  \end{scope}
\end{tikzpicture}

  \caption{ラウンドロビンの下での CPU バウンドのタスク T1 と I/O バウンド
    のタスク T2（\cref{ex:l07-io-bound}）．T2 がまだ I/O を待っている間に
    T1 のタイムスライスが満了した場合，T1 はそのまま実行を続ける．}
  \label{fig:l07-io-timeslice}
\end{figure}

\begin{example}
  \label{ex:l07-io-bound}
  CPU バウンドのタスク T1 が，\SI{5}{\milli\second} 計算してから
  \SI{10}{\milli\second} を要する I/O 操作を発行する I/O バウンドの
  タスク T2 と CPU を共有し，コンテキストスイッチには
  \SI{2}{\milli\second} を要するとする（\cref{fig:l07-io-timeslice}）．
  タイムスライスが \SI{5}{\milli\second} のとき，T2 の I/O 操作は T1 の
  2回目のタイムスライスの途中で完了し，T2 は再び実行されるまで
  \SI{4}{\milli\second} 待つ．このパターンは \SI{19}{\milli\second} ごとに
  繰り返される．T2 は \SI{19}{\milli\second} ごとに I/O 操作を発行し，
  デバイスは時間の $10/19 \approx \SI{53}{\percent}$ の間動いている．
  タイムスライスが \SI{30}{\milli\second} のときは，T2 は各 I/O 操作の
  後に \SI{24}{\milli\second} 待ち，パターンは \SI{39}{\milli\second} ごとに
  繰り返され，デバイスが動いているのは時間の
  $10/39 \approx \SI{26}{\percent}$ にすぎない．一方 T1 は，短いタイム
  スライスでは \SI{19}{\milli\second} ごとに \SI{10}{\milli\second}
  （$\approx\SI{53}{\percent}$）を得るが，長いタイムスライスでは
  \SI{39}{\milli\second} ごとに \SI{30}{\milli\second}
  （$\approx\SI{77}{\percent}$）を得る．
\end{example}

\subsection{まとめ}

2種類のタスクは正反対のタイムスライス長を好む．
\begin{description}
  \item[CPU バウンドのタスクは長いタイムスライスを好む．] コンテキスト
    スイッチが少ないほどオーバーヘッドが抑えられ，CPU 使用率とスループット
    が高く保たれる．
  \item[I/O バウンドのタスクは短いタイムスライスを好む．] I/O バウンドの
    タスクは I/O の完了後すぐに実行され，次の I/O 操作をより早く発行する．
    これによりデバイスが動き続け，対話的なタスクは典型的に I/O バウンドで
    あるから，利用者が体感する性能も良くなる．\margin{例の
    \SI{10}{\milli\second} はディスクの値である．\num{7200}~rpm の
    ドライブは平均しておよそ \SI{4}{\milli\second} の回転待ちと同程度の
    シークを要するが，NVMe の SSD は読み出しに数十マイクロ秒しか要しない．}
\end{description}

\subsection{CPU 使用率}

CPU 時間のうち有用な仕事に費やされる割合（\cref{eq:l02-useful}）は，
振る舞いの異なるタスクが混在する場合にも拡張できる．\tip{タイムスライス
$q$ ごとに費用 $c$ のスイッチが1回入るなら，この割合は $q/(q+c)$ である．
損失を \SI{1}{\percent} 未満に抑えるには，タイムスライスをコンテキスト
スイッチの費用の 100 倍以上にすればよい．}%
スケジュールの中で，
どのタスクも実行の機会を得ている繰り返しの区間を1つ選ぶ．その区間の中で，
CPU がタスクの実行に費やす総時間を $t_{\mathrm{run}}$，コンテキスト
スイッチに費やす総時間を $t_{\mathrm{ctx}}$ とすると，
\begin{equation}
  \text{CPU 使用率} \;=\;
  \frac{t_{\mathrm{run}}}{t_{\mathrm{run}} + t_{\mathrm{ctx}}}
  \label{eq:l07-utilization}
\end{equation}
である．

\begin{example}
  3つの I/O バウンドのタスクがそれぞれ \SI{2}{\milli\second} 計算して
  から \SI{5}{\milli\second} の I/O 操作を発行し，2つの CPU バウンドの
  タスクは決してブロックせず，コンテキストスイッチには
  \SI{0.5}{\milli\second} を要し，どのタスクも長時間実行されるとする．
  タイムスライスが \SI{2}{\milli\second} のときは，CPU バウンドかどうか
  によらずどのタスクも \SI{2}{\milli\second} で CPU を離れ，どの実行の
  後にもコンテキストスイッチが続くので，使用率は
  $2/(2+0.5) = \SI{80}{\percent}$ である．タイムスライスが
  \SI{20}{\milli\second} のときは，1巡は $3\cdot 2 + 2\cdot 20 =
  \SI{46}{\milli\second}$ の計算と5回のコンテキストスイッチからなり，
  使用率は $46/(46+2.5) \approx \SI{94.8}{\percent}$ である．12.5 と
  \SI{48.5}{\milli\second} というどちらの1巡も I/O 操作より長いので，
  I/O バウンドのタスクは自分の順番が来たときには常に実行可能である．%
  \caution{\cref{eq:l07-utilization} はアイドル時間を無視している．どの
  タスクも I/O を待っている区間は $t_{\mathrm{run}}$ にも
  $t_{\mathrm{ctx}}$ にも寄与しないので，CPU に実行するものがない時間が
  長くても，この比は \SI{100}{\percent} になりうる．}
\end{example}

\begin{keypoint}
  タイムシェアリングは，実行時間の知識なしにシステムを応答性の高いものに
  するが，その代償としてタイムスライスごとに割り込み，スケジューラの実行，
  コンテキストスイッチが必要になる．CPU バウンドのタスクは長いタイム
  スライスを，I/O バウンドのタスクは短いタイムスライスを好む．
\end{keypoint}

\needspace{7\baselineskip}
\section{多段フィードバックキュー}
\label{sec:l07-mlfq}

I/O バウンドのタスクと CPU バウンドのタスクに異なるタイムスライスを与える
には，スケジューラはすべてのタスクを1つのランキューに保持して選択時に
各タスクの種類を調べてもよいし，2種類のタスクを別々の構造に保持しても
よい．\source{P3L1，ランキューのデータ構造；OSTEP ch.~8；Corbató ら
\cite{corbato1962}．}複数のキューからなるランキューは後者の方法を取る
（\cref{fig:l07-mlfq}）．
\begin{itemize}
  \item 1番目のキューは最も I/O 集約的なタスクを保持し，最も高い優先度を
    持ち，そのタスクに \SI{8}{\milli\second} のタイムスライスを与える．
  \item 2番目のキューは I/O と計算が混在するタスクを保持し，タイム
    スライスは \SI{16}{\milli\second} である．
  \item 最後のキューは CPU 集約的なタスクを保持し，最も低い優先度を持ち，
    無限のタイムスライス，すなわち FCFS でそれらをスケジュールする．
\end{itemize}
I/O バウンドのタスクはタイムシェアリングの利点を受け，CPU バウンドの
タスクはそのオーバーヘッドを避けられる．

\begin{figure}[htbp]
  \centering
  % Multi-level feedback queue with three levels.
% Usage: % Multi-level feedback queue with three levels.
% Usage: % Multi-level feedback queue with three levels.
% Usage: \input{figures/l07-mlfq}   (width ~104 mm)
\begin{tikzpicture}[x=1mm, y=1mm, font=\footnotesize\sffamily,
  >={Stealth[length=1.6mm]},
  desc/.style={anchor=east, align=right, font=\scriptsize\sffamily},
  q/.style={draw, fill=ln-box},
  tk/.style={draw, circle, fill=white, minimum size=3.2mm, inner sep=0pt}]
  \foreach \y in {28,14,0} {
    \draw[q] (0,\y-3) rectangle (36,\y+3);
    \foreach \i in {1,...,4} \draw (36-5*\i,\y-3) -- ++(0,6);
    \draw[->] (36,\y) -- (44,\y);
  }
  \node[tk] at (33.5,28) {}; \node[tk] at (28.5,28) {}; \node[tk] at (23.5,28) {};
  \node[tk] at (33.5,14) {}; \node[tk] at (28.5,14) {};
  \node[tk] at (33.5,0) {};
  \node[anchor=west] at (44.5,28) {CPU};
  \node[anchor=west] at (44.5,14) {CPU};
  \node[anchor=west] at (44.5,0) {CPU};
  \node[desc] at (-3,28) {\iflangja{タイムスライス 8\,ms\\最も I/O 集約的\\最高優先度}{timeslice 8\,ms\\most I/O-intensive\\highest priority}};
  \node[desc] at (-3,14) {\iflangja{タイムスライス 16\,ms\\I/O と CPU の混合}{timeslice 16\,ms\\mix of I/O and CPU}};
  \node[desc] at (-3,0)  {\iflangja{タイムスライス $\infty$（FCFS）\\CPU 集約的\\最低優先度}{timeslice $\infty$ (FCFS)\\CPU-intensive\\lowest priority}};
  % new tasks enter the top level
  \draw[->] (2,40) node[above, font=\scriptsize\sffamily] {\iflangja{新しいタスク}{new task}} -- (2,31.5);
  % feedback
  \draw[->, thick] (56,27) to[bend left=40] (56,15);
  \draw[->, thick] (56,13) to[bend left=40] (56,1);
  \draw[->, thick, densely dashed] (66,1) to[bend right=40] (66,13);
  \draw[->, thick, densely dashed] (66,15) to[bend right=40] (66,27);
  % legend
  \begin{scope}[yshift=-12mm, font=\scriptsize\sffamily]
    \draw[->, thick] (0,0) -- (7,0);
    \node[anchor=west] at (8,0) {\iflangja{タイムスライスを使い切る：下の段へ}{uses its whole timeslice: moves down}};
    \draw[->, thick, densely dashed] (0,-5) -- (7,-5);
    \node[anchor=west] at (8,-5) {\iflangja{I/O のため早く CPU を手放す：上の段へ}{releases the CPU early for I/O: moves up}};
  \end{scope}
\end{tikzpicture}
   (width ~104 mm)
\begin{tikzpicture}[x=1mm, y=1mm, font=\footnotesize\sffamily,
  >={Stealth[length=1.6mm]},
  desc/.style={anchor=east, align=right, font=\scriptsize\sffamily},
  q/.style={draw, fill=ln-box},
  tk/.style={draw, circle, fill=white, minimum size=3.2mm, inner sep=0pt}]
  \foreach \y in {28,14,0} {
    \draw[q] (0,\y-3) rectangle (36,\y+3);
    \foreach \i in {1,...,4} \draw (36-5*\i,\y-3) -- ++(0,6);
    \draw[->] (36,\y) -- (44,\y);
  }
  \node[tk] at (33.5,28) {}; \node[tk] at (28.5,28) {}; \node[tk] at (23.5,28) {};
  \node[tk] at (33.5,14) {}; \node[tk] at (28.5,14) {};
  \node[tk] at (33.5,0) {};
  \node[anchor=west] at (44.5,28) {CPU};
  \node[anchor=west] at (44.5,14) {CPU};
  \node[anchor=west] at (44.5,0) {CPU};
  \node[desc] at (-3,28) {\iflangja{タイムスライス 8\,ms\\最も I/O 集約的\\最高優先度}{timeslice 8\,ms\\most I/O-intensive\\highest priority}};
  \node[desc] at (-3,14) {\iflangja{タイムスライス 16\,ms\\I/O と CPU の混合}{timeslice 16\,ms\\mix of I/O and CPU}};
  \node[desc] at (-3,0)  {\iflangja{タイムスライス $\infty$（FCFS）\\CPU 集約的\\最低優先度}{timeslice $\infty$ (FCFS)\\CPU-intensive\\lowest priority}};
  % new tasks enter the top level
  \draw[->] (2,40) node[above, font=\scriptsize\sffamily] {\iflangja{新しいタスク}{new task}} -- (2,31.5);
  % feedback
  \draw[->, thick] (56,27) to[bend left=40] (56,15);
  \draw[->, thick] (56,13) to[bend left=40] (56,1);
  \draw[->, thick, densely dashed] (66,1) to[bend right=40] (66,13);
  \draw[->, thick, densely dashed] (66,15) to[bend right=40] (66,27);
  % legend
  \begin{scope}[yshift=-12mm, font=\scriptsize\sffamily]
    \draw[->, thick] (0,0) -- (7,0);
    \node[anchor=west] at (8,0) {\iflangja{タイムスライスを使い切る：下の段へ}{uses its whole timeslice: moves down}};
    \draw[->, thick, densely dashed] (0,-5) -- (7,-5);
    \node[anchor=west] at (8,-5) {\iflangja{I/O のため早く CPU を手放す：上の段へ}{releases the CPU early for I/O: moves up}};
  \end{scope}
\end{tikzpicture}
   (width ~104 mm)
\begin{tikzpicture}[x=1mm, y=1mm, font=\footnotesize\sffamily,
  >={Stealth[length=1.6mm]},
  desc/.style={anchor=east, align=right, font=\scriptsize\sffamily},
  q/.style={draw, fill=ln-box},
  tk/.style={draw, circle, fill=white, minimum size=3.2mm, inner sep=0pt}]
  \foreach \y in {28,14,0} {
    \draw[q] (0,\y-3) rectangle (36,\y+3);
    \foreach \i in {1,...,4} \draw (36-5*\i,\y-3) -- ++(0,6);
    \draw[->] (36,\y) -- (44,\y);
  }
  \node[tk] at (33.5,28) {}; \node[tk] at (28.5,28) {}; \node[tk] at (23.5,28) {};
  \node[tk] at (33.5,14) {}; \node[tk] at (28.5,14) {};
  \node[tk] at (33.5,0) {};
  \node[anchor=west] at (44.5,28) {CPU};
  \node[anchor=west] at (44.5,14) {CPU};
  \node[anchor=west] at (44.5,0) {CPU};
  \node[desc] at (-3,28) {\iflangja{タイムスライス 8\,ms\\最も I/O 集約的\\最高優先度}{timeslice 8\,ms\\most I/O-intensive\\highest priority}};
  \node[desc] at (-3,14) {\iflangja{タイムスライス 16\,ms\\I/O と CPU の混合}{timeslice 16\,ms\\mix of I/O and CPU}};
  \node[desc] at (-3,0)  {\iflangja{タイムスライス $\infty$（FCFS）\\CPU 集約的\\最低優先度}{timeslice $\infty$ (FCFS)\\CPU-intensive\\lowest priority}};
  % new tasks enter the top level
  \draw[->] (2,40) node[above, font=\scriptsize\sffamily] {\iflangja{新しいタスク}{new task}} -- (2,31.5);
  % feedback
  \draw[->, thick] (56,27) to[bend left=40] (56,15);
  \draw[->, thick] (56,13) to[bend left=40] (56,1);
  \draw[->, thick, densely dashed] (66,1) to[bend right=40] (66,13);
  \draw[->, thick, densely dashed] (66,15) to[bend right=40] (66,27);
  % legend
  \begin{scope}[yshift=-12mm, font=\scriptsize\sffamily]
    \draw[->, thick] (0,0) -- (7,0);
    \node[anchor=west] at (8,0) {\iflangja{タイムスライスを使い切る：下の段へ}{uses its whole timeslice: moves down}};
    \draw[->, thick, densely dashed] (0,-5) -- (7,-5);
    \node[anchor=west] at (8,-5) {\iflangja{I/O のため早く CPU を手放す：上の段へ}{releases the CPU early for I/O: moves up}};
  \end{scope}
\end{tikzpicture}

  \caption{3段の多段フィードバックキュー．各段はそれぞれのタイムスライス
    でタスクを処理し，タスクはタイムスライスの使い方に応じて段の間を
    移動する．}
  \label{fig:l07-mlfq}
\end{figure}

この設計は，あるタスクが I/O 集約的か CPU 集約的かをスケジューラがどう
知るのか，という問題を提起する．タスクの履歴に基づく経験則は，新しい
タスクにも，振る舞いを変えるタスクにも役立たない．その解決策は，タスクが
タイムスライスをどう使うかから，スケジューラにタスクの種類を学習させる
ことである．\caution{OSTEP の MLFQ は2点で異なる．どの段もタスクを
ラウンドロビンで実行し，最下段も例外ではない．また周期 $S$ ごとにすべての
タスクが最上段へ戻され，どのタスクもスタベーションに陥らない（OSTEP
ch.~8）．}
\begin{enumerate}
  \item 新しいタスクは最上段のキューに入る．
  \item タスクがタイムスライスの満了前に CPU を手放したなら，段の選択は
    適切であったのであり，そのタスクはこの段に留まる．
  \item タスクがタイムスライス全体を使い切ったなら，それはその段が想定
    するより CPU 集約的であり，1段下へ移される．
  \item 下の段にあるタスクが I/O を待つために早く CPU を解放したなら，
    優先度の引き上げを受けて上へ移される．
\end{enumerate}

\begin{definition}[多段フィードバックキュー]
  \term[ただんふぃーどばっくきゅー]{多段フィードバックキュー}%
  [multi-level feedback queue]%
  \index{MLFQ|see{多段フィードバックキュー}}（MLFQ）とは，複数の段から
  なり，各段が独自のスケジューリング方針を持ち，観測された CPU の使い方に
  応じてタスクが段の間を移動するランキューである．
\end{definition}

MLFQ は単なる優先度キューの集まりではない．その各段は，長さの異なる
タイムスライスと最下段の FCFS によってタスクを異なる仕方で扱い，
フィードバック機構が時間とともに各タスクの段を調節する．この設計は
Fernando Corbató のタイムシェアリングシステム CTSS \cite{corbato1962} に
さかのぼる．
Corbató は 1990 年にチューリング賞を受賞した．

\begin{example}
  新しいタスクが \SI{8}{\milli\second} の段に入り，ブロックせずに
  \SI{30}{\milli\second} 計算するとする．このタスクは最初のタイムスライス
  全体を使い切って \SI{16}{\milli\second} の段へ移り，そのタイムスライスも
  使い切って最下段へ移り，そこで残りの \SI{6}{\milli\second} を完了する．
  I/O 操作の間に \SI{3}{\milli\second} 計算するタスクは，
  \SI{8}{\milli\second} のタイムスライスを使い切ることが決してなく，最上段
  に留まる．
\end{example}

\needspace{7\baselineskip}
\section{Linux の $O(1)$ スケジューラ}
\label{sec:l07-o1}

Linux の \term[O1すけじゅーら]{$O(1)$ スケジューラ}[O(1) scheduler]は，
システム内のタスク数によらず，実行するタスクの選択とランキューへのタスク
の追加を定数時間で行う．\source{P3L1，Linux の O(1) スケジューラ；Bovet
と Cesati \cite{bovet2005} ch.~7．}これはプリエンプティブで優先度に基づく
スケジューラであり，140 段階の優先度を持ち，番号が小さいほど優先度が高い
ように番号付けられている．0 から 99 までの水準はリアルタイムタスク用，
100 から 139 までの水準はタイムシェアリングのタスク用である．

ユーザプロセスの既定の優先度は 120 である．これはプロセスの
\term[niceち]{nice値}[nice value]によって調整でき，nice 値は $-20$ から
19 の範囲を取り，既定値に加算されるので，タイムシェアリングのタスクの
優先度は 100 から 139 の範囲となる．プロセスは，例えば \texttt{nice}
システムコールによって自分の優先度を下げられるが，上げるには権限が必要で
ある．

タスクのタイムスライスはその優先度に依存する．優先度の低いタスクには最も
短いタイムスライスが，優先度の高いタスクには最も長いタイムスライスが与え
られる．スケジューラはさらに，タスクの\emph{スリープ時間}，すなわちタスク
が待機またはアイドルで過ごす時間に基づくフィードバックも用いる．
\begin{itemize}
  \item 長い時間スリープするタスクは対話的であり，その優先度は最大 5 水準
    引き上げられる（その番号は最大 5 減る）．
  \item ほとんどスリープしないタスクは計算集約的であり，その優先度は最大
    5 水準引き下げられる（その番号は最大 5 増える）．
\end{itemize}

\begin{example}
  Bovet と Cesati が記述している Linux~2.6 カーネルでは，nice 値が与える
  優先度である静的優先度が $p$ であるタイムシェアリングのタスクには，
  $p < 120$ なら $(140 - p) \times \SI{20}{\milli\second}$ の，そうで
  なければ $(140 - p) \times \SI{5}{\milli\second}$ のタイムスライスが
  与えられる．nice 値 5 のプロセスは優先度 125 を持ち，タイムスライスは
  $15 \times 5 = \SI{75}{\milli\second}$ である．nice 値 0 のプロセスは
  \SI{100}{\milli\second}，nice 値 $-10$ のプロセスは
  $30 \times 20 = \SI{600}{\milli\second}$ である．nice 値 5 のプロセスが
  対話的であれば，フィードバックによって優先度は 120 まで高められ，
  スリープせずに計算するのであれば 130 まで下げられる．\caution{2つの
  番号付けは向きが逆である．優先度の番号が大きいほど，すなわち nice 値が
  大きいほど優先度は低いが，\texttt{SCHED\_FIFO} と \texttt{SCHED\_RR}
  が用いる \texttt{sched\_priority}（1--99）は大きいほど優先度が高い．}
\end{example}

\begin{figure}[htbp]
  \centering
  % Runqueue of the Linux O(1) scheduler: an active and an expired array of
% 140 priority lists, each with a bitmap of the non-empty lists.
% Usage: % Runqueue of the Linux O(1) scheduler: an active and an expired array of
% 140 priority lists, each with a bitmap of the non-empty lists.
% Usage: % Runqueue of the Linux O(1) scheduler: an active and an expired array of
% 140 priority lists, each with a bitmap of the non-empty lists.
% Usage: \input{figures/l07-o1}   (width ~112 mm)
\begin{tikzpicture}[x=1mm, y=1mm, font=\scriptsize\sffamily,
  >={Stealth[length=1.5mm]},
  bit/.style={draw, minimum width=4mm, minimum height=4mm, inner sep=0pt},
  slot/.style={draw, fill=ln-box, minimum width=5mm, minimum height=4mm, inner sep=0pt},
  tk/.style={draw, circle, fill=ln-accent!22, minimum size=3.6mm, inner sep=0pt},
  ttl/.style={font=\footnotesize\sffamily\bfseries}]
  % rows: label / y
  \def\rows{0/0, {\vdots}/-4.5, 99/-9, 100/-13.5, {\vdots}/-18, 120/-22.5, {\vdots}/-27, 139/-31.5}
  % ---------------- active array at x0 = 18
  \node[ttl, anchor=south] at (32,5) {\iflangja{アクティブ配列}{active array}};
  \node[anchor=south, text=ln-muted] at (18,2) {\iflangja{ビットマップ}{bitmap}};
  \foreach \l/\y in \rows { \node[anchor=east] at (29,\y) {\l}; }
  \foreach \y/\b in {0/0, -9/0, -13.5/1, -22.5/1, -31.5/1} {
    \node[bit] at (18,\y) {\b};
    \node[slot] at (33,\y) {};
  }
  \node[tk] (a1) at (41,-13.5) {};
  \draw[->] (33,-13.5) -- (a1);
  \node[tk] (a2) at (41,-22.5) {};
  \node[tk] (a3) at (48,-22.5) {};
  \draw[->] (33,-22.5) -- (a2); \draw[->] (a2) -- (a3);
  \node[tk] (a4) at (41,-31.5) {};
  \draw[->] (33,-31.5) -- (a4);
  % priority classes
  \draw (13,2) -- (11.5,2) -- (11.5,-11) -- (13,-11);
  \node[anchor=east, align=right] at (9.5,-4.5) {\iflangja{リアル\\タイム}{real-\\time}};
  \draw (13,-11.8) -- (11.5,-11.8) -- (11.5,-33.5) -- (13,-33.5);
  \node[anchor=east, align=right] at (9.5,-22.5) {\iflangja{タイム\\シェアリング}{time-\\sharing}};
  % ---------------- expired array at x0 = 76
  \node[ttl, anchor=south] at (90,5) {\iflangja{期限切れ配列}{expired array}};
  \node[anchor=south, text=ln-muted] at (76,2) {\iflangja{ビットマップ}{bitmap}};
  \foreach \l/\y in \rows { \node[anchor=east] at (87,\y) {\l}; }
  \foreach \y/\b in {0/0, -9/0, -13.5/0, -22.5/1, -31.5/0} {
    \node[bit] at (76,\y) {\b};
    \node[slot] at (91,\y) {};
  }
  \node[tk] (e1) at (99,-22.5) {};
  \draw[->] (91,-22.5) -- (e1);
  % a task whose timeslice expires moves to the expired array
  \draw[->, thick, ln-accent, rounded corners=2mm] (a3.south) -- (48,-37.5) -- (99,-37.5) -- (e1.south);
  \node[above, text=ln-accent, inner sep=1pt] at (62,-37.5) {\iflangja{タイムスライス満了}{timeslice expires}};
  % swap
  \draw[<->, thick] (44,-44) -- node[below, align=center] {\iflangja{アクティブ配列が空になったら交換}{swapped when the active array is empty}} (88,-44);
\end{tikzpicture}
   (width ~112 mm)
\begin{tikzpicture}[x=1mm, y=1mm, font=\scriptsize\sffamily,
  >={Stealth[length=1.5mm]},
  bit/.style={draw, minimum width=4mm, minimum height=4mm, inner sep=0pt},
  slot/.style={draw, fill=ln-box, minimum width=5mm, minimum height=4mm, inner sep=0pt},
  tk/.style={draw, circle, fill=ln-accent!22, minimum size=3.6mm, inner sep=0pt},
  ttl/.style={font=\footnotesize\sffamily\bfseries}]
  % rows: label / y
  \def\rows{0/0, {\vdots}/-4.5, 99/-9, 100/-13.5, {\vdots}/-18, 120/-22.5, {\vdots}/-27, 139/-31.5}
  % ---------------- active array at x0 = 18
  \node[ttl, anchor=south] at (32,5) {\iflangja{アクティブ配列}{active array}};
  \node[anchor=south, text=ln-muted] at (18,2) {\iflangja{ビットマップ}{bitmap}};
  \foreach \l/\y in \rows { \node[anchor=east] at (29,\y) {\l}; }
  \foreach \y/\b in {0/0, -9/0, -13.5/1, -22.5/1, -31.5/1} {
    \node[bit] at (18,\y) {\b};
    \node[slot] at (33,\y) {};
  }
  \node[tk] (a1) at (41,-13.5) {};
  \draw[->] (33,-13.5) -- (a1);
  \node[tk] (a2) at (41,-22.5) {};
  \node[tk] (a3) at (48,-22.5) {};
  \draw[->] (33,-22.5) -- (a2); \draw[->] (a2) -- (a3);
  \node[tk] (a4) at (41,-31.5) {};
  \draw[->] (33,-31.5) -- (a4);
  % priority classes
  \draw (13,2) -- (11.5,2) -- (11.5,-11) -- (13,-11);
  \node[anchor=east, align=right] at (9.5,-4.5) {\iflangja{リアル\\タイム}{real-\\time}};
  \draw (13,-11.8) -- (11.5,-11.8) -- (11.5,-33.5) -- (13,-33.5);
  \node[anchor=east, align=right] at (9.5,-22.5) {\iflangja{タイム\\シェアリング}{time-\\sharing}};
  % ---------------- expired array at x0 = 76
  \node[ttl, anchor=south] at (90,5) {\iflangja{期限切れ配列}{expired array}};
  \node[anchor=south, text=ln-muted] at (76,2) {\iflangja{ビットマップ}{bitmap}};
  \foreach \l/\y in \rows { \node[anchor=east] at (87,\y) {\l}; }
  \foreach \y/\b in {0/0, -9/0, -13.5/0, -22.5/1, -31.5/0} {
    \node[bit] at (76,\y) {\b};
    \node[slot] at (91,\y) {};
  }
  \node[tk] (e1) at (99,-22.5) {};
  \draw[->] (91,-22.5) -- (e1);
  % a task whose timeslice expires moves to the expired array
  \draw[->, thick, ln-accent, rounded corners=2mm] (a3.south) -- (48,-37.5) -- (99,-37.5) -- (e1.south);
  \node[above, text=ln-accent, inner sep=1pt] at (62,-37.5) {\iflangja{タイムスライス満了}{timeslice expires}};
  % swap
  \draw[<->, thick] (44,-44) -- node[below, align=center] {\iflangja{アクティブ配列が空になったら交換}{swapped when the active array is empty}} (88,-44);
\end{tikzpicture}
   (width ~112 mm)
\begin{tikzpicture}[x=1mm, y=1mm, font=\scriptsize\sffamily,
  >={Stealth[length=1.5mm]},
  bit/.style={draw, minimum width=4mm, minimum height=4mm, inner sep=0pt},
  slot/.style={draw, fill=ln-box, minimum width=5mm, minimum height=4mm, inner sep=0pt},
  tk/.style={draw, circle, fill=ln-accent!22, minimum size=3.6mm, inner sep=0pt},
  ttl/.style={font=\footnotesize\sffamily\bfseries}]
  % rows: label / y
  \def\rows{0/0, {\vdots}/-4.5, 99/-9, 100/-13.5, {\vdots}/-18, 120/-22.5, {\vdots}/-27, 139/-31.5}
  % ---------------- active array at x0 = 18
  \node[ttl, anchor=south] at (32,5) {\iflangja{アクティブ配列}{active array}};
  \node[anchor=south, text=ln-muted] at (18,2) {\iflangja{ビットマップ}{bitmap}};
  \foreach \l/\y in \rows { \node[anchor=east] at (29,\y) {\l}; }
  \foreach \y/\b in {0/0, -9/0, -13.5/1, -22.5/1, -31.5/1} {
    \node[bit] at (18,\y) {\b};
    \node[slot] at (33,\y) {};
  }
  \node[tk] (a1) at (41,-13.5) {};
  \draw[->] (33,-13.5) -- (a1);
  \node[tk] (a2) at (41,-22.5) {};
  \node[tk] (a3) at (48,-22.5) {};
  \draw[->] (33,-22.5) -- (a2); \draw[->] (a2) -- (a3);
  \node[tk] (a4) at (41,-31.5) {};
  \draw[->] (33,-31.5) -- (a4);
  % priority classes
  \draw (13,2) -- (11.5,2) -- (11.5,-11) -- (13,-11);
  \node[anchor=east, align=right] at (9.5,-4.5) {\iflangja{リアル\\タイム}{real-\\time}};
  \draw (13,-11.8) -- (11.5,-11.8) -- (11.5,-33.5) -- (13,-33.5);
  \node[anchor=east, align=right] at (9.5,-22.5) {\iflangja{タイム\\シェアリング}{time-\\sharing}};
  % ---------------- expired array at x0 = 76
  \node[ttl, anchor=south] at (90,5) {\iflangja{期限切れ配列}{expired array}};
  \node[anchor=south, text=ln-muted] at (76,2) {\iflangja{ビットマップ}{bitmap}};
  \foreach \l/\y in \rows { \node[anchor=east] at (87,\y) {\l}; }
  \foreach \y/\b in {0/0, -9/0, -13.5/0, -22.5/1, -31.5/0} {
    \node[bit] at (76,\y) {\b};
    \node[slot] at (91,\y) {};
  }
  \node[tk] (e1) at (99,-22.5) {};
  \draw[->] (91,-22.5) -- (e1);
  % a task whose timeslice expires moves to the expired array
  \draw[->, thick, ln-accent, rounded corners=2mm] (a3.south) -- (48,-37.5) -- (99,-37.5) -- (e1.south);
  \node[above, text=ln-accent, inner sep=1pt] at (62,-37.5) {\iflangja{タイムスライス満了}{timeslice expires}};
  % swap
  \draw[<->, thick] (44,-44) -- node[below, align=center] {\iflangja{アクティブ配列が空になったら交換}{swapped when the active array is empty}} (88,-44);
\end{tikzpicture}

  \caption{$O(1)$ スケジューラのランキュー．各配列は優先度の水準ごとに
    1本のリストと，空でないリストを示すビットマップを持つ．タイムスライス
    が満了したタスクは期限切れ配列へ移り，アクティブ配列が空になると2つの
    配列が交換される．}
  \label{fig:l07-o1}
\end{figure}

ランキューは，優先度の水準ごとに1本のリストを持つ，タスクのリストの配列
2つからなる（\cref{fig:l07-o1}）．スケジューラが次のタスクを取り出す方の
配列を\term[あくてぃぶはいれつ]{アクティブ配列}[active array]という．
ビットマップがその 140 本のリストのうちどれが空でないかを記録しており，
スケジューラは最初に立っているビットを見つけてそのリストの先頭のタスクを
取り出す．タスクの追加はリストの末尾への追加とビットの設定である．水準の
数が固定であるから，どちらの操作も定数時間で済む．タスクは，そのタイム
スライスが満了するまでアクティブ配列に留まる．それより前にブロックした
タスクは，再び実行可能になったときに，残りのタイムスライスとともに
アクティブ配列へ戻る．

\term[きげんぎれはいれつ]{期限切れ配列}[expired array]は，タイムスライス
が満了して非アクティブとなったタスクを保持する．タスクのタイムスライスが
満了すると，スケジューラはその優先度とタイムスライスを再計算し，タスクを
期限切れ配列へ移す．\margin{Linux~2.6 カーネルでは，対話的と判定された
タスクは，期限切れ配列のタスクがすでに待ちすぎている場合を除き，通常は
アクティブ配列へ戻された．}アクティブ配列にタスクが残らなくなると，
スケジューラは2つの配列を交換する．これにはポインタを2つ入れ替えるだけで
よい．この仕組みは，優先度の低いタスクにも実行の機会を与える．優先度の
高いタスクは長いタイムスライスを受け取るものの，それを使い切れば期限切れ
配列へ移り，アクティブ配列のすべてのタスクが，優先度の低いものも含めて，
自分のタイムスライスを使い切るまでそこで待つからである．

$O(1)$ スケジューラは Ingo Molnar によって Linux~2.5 の開発系列で導入され，
Linux~2.6 のスケジューラとなった．ストリーミングメディアのような
アプリケーションによってワークロードがより対話的になるにつれ，その
ジッタが問題となり，Linux~2.6.23 では，同じく Molnar が書いた Completely
Fair Scheduler に置き換えられた．

\needspace{7\baselineskip}
\section{Linux の Completely Fair Scheduler}
\label{sec:l07-cfs}

$O(1)$ スケジューラには2つの問題がある．\source{P3L1，Linux の CFS
スケジューラ；OSTEP ch.~9．}第1に，対話的なタスクの性能が悪い．いったん
タスクが期限切れ配列に入ると，アクティブ配列のすべてのタスクがタイム
スライスを使い切るまでスケジュールされないので，対話的なタスクが実行され
るまでに予測できない時間待つことがある．第2に，公平性の保証がない．ここ
で公平性とは，与えられた時間区間の中で，どのタスクもその優先度に比例した
時間だけ実行されるべきだということである．

\term{CFS}[Completely Fair Scheduler]%
\index{Completely Fair Scheduler|see{CFS}}（完全公平スケジューラ）は，
その両方の問題に対処する．これはバージョン 2.6.23 から 6.5 まで，
リアルタイムでないタスクのための Linux の既定のスケジューラであった．その
ランキューは\term[あかくろぎ]{赤黒木}[red-black tree]である．これは自己
平衡二分探索木であり，根から葉へのどの経路も他のどの経路の2倍より長く
ならないので，すべての経路の長さがほぼ等しくなる．木はタスクの仮想実行
時間で順序付けられている（\cref{fig:l07-cfs}）．

\begin{definition}[仮想実行時間]
  タスクの\term[かそうじっこうじかん]{仮想実行時間}[virtual runtime]%
  （\texttt{vruntime}）とは，そのタスクが CPU 上で実行に費やした時間で
  あり，ナノ秒の粒度で測られ，タスクの優先度に応じて重み付けされる．
\end{definition}

\begin{figure}[htbp]
  \centering
  % CFS runqueue: a red-black tree of the ready tasks ordered by virtual
% runtime (ms).
% Usage: % CFS runqueue: a red-black tree of the ready tasks ordered by virtual
% runtime (ms).
% Usage: % CFS runqueue: a red-black tree of the ready tasks ordered by virtual
% runtime (ms).
% Usage: \input{figures/l07-cfs}   (width ~108 mm)
\begin{tikzpicture}[x=1mm, y=1mm, font=\footnotesize\sffamily,
  >={Stealth[length=1.6mm]},
  bn/.style={draw, circle, fill=ln-muted, text=white, minimum size=7mm, inner sep=0pt},
  rn/.style={draw, circle, fill=ln-caution!80, text=white, minimum size=7mm, inner sep=0pt},
  lab/.style={font=\scriptsize\sffamily, align=center}]
  \node[bn] (n24) at (48,0)   {24};
  \node[bn] (n15) at (28,-13) {15};
  \node[bn] (n38) at (68,-13) {38};
  \node[rn] (n9)  at (18,-26) {9};
  \node[rn] (n19) at (38,-26) {19};
  \node[rn] (n30) at (58,-26) {30};
  \node[rn] (n45) at (78,-26) {45};
  \draw (n24) -- (n15); \draw (n24) -- (n38);
  \draw (n15) -- (n9);  \draw (n15) -- (n19);
  \draw (n38) -- (n30); \draw (n38) -- (n45);
  % leftmost node
  \draw[->, thick, ln-accent] (n9.west) -- ++(-9,0) node[left, lab, text=ln-accent]
    {\iflangja{最左の節点：\\vruntime 最小\\→ 次に実行}{leftmost node:\\smallest vruntime\\$\to$ runs next}};
  % a preempted task, whose vruntime has just passed that of the leftmost
  % node, is reinserted
  \node[bn, dashed, fill=white, text=black, draw=black] (p) at (98,-8) {12};
  \node[lab, anchor=south] at (98,-3.5) {\iflangja{プリエンプト\\されたタスク}{preempted\\task}};
  \node[draw, dashed, circle, minimum size=7mm, inner sep=0pt] (ins) at (25,-38) {};
  \draw[dashed] (n9) -- (ins);
  \draw[->, thick, densely dashed] (p.south) .. controls (98,-48) and (58,-48) .. (ins.east);
  \node[lab] at (72,-38) {\iflangja{挿入}{insert}};
  \node[lab, text=ln-muted, anchor=west] at (0,-48) {\iflangja{節点の値：仮想実行時間（ms）}{node values: virtual runtime (ms)}};
\end{tikzpicture}
   (width ~108 mm)
\begin{tikzpicture}[x=1mm, y=1mm, font=\footnotesize\sffamily,
  >={Stealth[length=1.6mm]},
  bn/.style={draw, circle, fill=ln-muted, text=white, minimum size=7mm, inner sep=0pt},
  rn/.style={draw, circle, fill=ln-caution!80, text=white, minimum size=7mm, inner sep=0pt},
  lab/.style={font=\scriptsize\sffamily, align=center}]
  \node[bn] (n24) at (48,0)   {24};
  \node[bn] (n15) at (28,-13) {15};
  \node[bn] (n38) at (68,-13) {38};
  \node[rn] (n9)  at (18,-26) {9};
  \node[rn] (n19) at (38,-26) {19};
  \node[rn] (n30) at (58,-26) {30};
  \node[rn] (n45) at (78,-26) {45};
  \draw (n24) -- (n15); \draw (n24) -- (n38);
  \draw (n15) -- (n9);  \draw (n15) -- (n19);
  \draw (n38) -- (n30); \draw (n38) -- (n45);
  % leftmost node
  \draw[->, thick, ln-accent] (n9.west) -- ++(-9,0) node[left, lab, text=ln-accent]
    {\iflangja{最左の節点：\\vruntime 最小\\→ 次に実行}{leftmost node:\\smallest vruntime\\$\to$ runs next}};
  % a preempted task, whose vruntime has just passed that of the leftmost
  % node, is reinserted
  \node[bn, dashed, fill=white, text=black, draw=black] (p) at (98,-8) {12};
  \node[lab, anchor=south] at (98,-3.5) {\iflangja{プリエンプト\\されたタスク}{preempted\\task}};
  \node[draw, dashed, circle, minimum size=7mm, inner sep=0pt] (ins) at (25,-38) {};
  \draw[dashed] (n9) -- (ins);
  \draw[->, thick, densely dashed] (p.south) .. controls (98,-48) and (58,-48) .. (ins.east);
  \node[lab] at (72,-38) {\iflangja{挿入}{insert}};
  \node[lab, text=ln-muted, anchor=west] at (0,-48) {\iflangja{節点の値：仮想実行時間（ms）}{node values: virtual runtime (ms)}};
\end{tikzpicture}
   (width ~108 mm)
\begin{tikzpicture}[x=1mm, y=1mm, font=\footnotesize\sffamily,
  >={Stealth[length=1.6mm]},
  bn/.style={draw, circle, fill=ln-muted, text=white, minimum size=7mm, inner sep=0pt},
  rn/.style={draw, circle, fill=ln-caution!80, text=white, minimum size=7mm, inner sep=0pt},
  lab/.style={font=\scriptsize\sffamily, align=center}]
  \node[bn] (n24) at (48,0)   {24};
  \node[bn] (n15) at (28,-13) {15};
  \node[bn] (n38) at (68,-13) {38};
  \node[rn] (n9)  at (18,-26) {9};
  \node[rn] (n19) at (38,-26) {19};
  \node[rn] (n30) at (58,-26) {30};
  \node[rn] (n45) at (78,-26) {45};
  \draw (n24) -- (n15); \draw (n24) -- (n38);
  \draw (n15) -- (n9);  \draw (n15) -- (n19);
  \draw (n38) -- (n30); \draw (n38) -- (n45);
  % leftmost node
  \draw[->, thick, ln-accent] (n9.west) -- ++(-9,0) node[left, lab, text=ln-accent]
    {\iflangja{最左の節点：\\vruntime 最小\\→ 次に実行}{leftmost node:\\smallest vruntime\\$\to$ runs next}};
  % a preempted task, whose vruntime has just passed that of the leftmost
  % node, is reinserted
  \node[bn, dashed, fill=white, text=black, draw=black] (p) at (98,-8) {12};
  \node[lab, anchor=south] at (98,-3.5) {\iflangja{プリエンプト\\されたタスク}{preempted\\task}};
  \node[draw, dashed, circle, minimum size=7mm, inner sep=0pt] (ins) at (25,-38) {};
  \draw[dashed] (n9) -- (ins);
  \draw[->, thick, densely dashed] (p.south) .. controls (98,-48) and (58,-48) .. (ins.east);
  \node[lab] at (72,-38) {\iflangja{挿入}{insert}};
  \node[lab, text=ln-muted, anchor=west] at (0,-48) {\iflangja{節点の値：仮想実行時間（ms）}{node values: virtual runtime (ms)}};
\end{tikzpicture}

  \caption{CFS のランキュー．実行可能なタスクを仮想実行時間で順序付けた
    赤黒木である（灰色の節点が黒，赤い節点が赤）．最左のタスクが次に実行
    され，プリエンプトされたタスクはその仮想実行時間が定める位置に挿入
    される．}
  \label{fig:l07-cfs}
\end{figure}

CFS のアルゴリズムは単純である．\margin{Linux~6.6 以降，公平スケジュー
リングクラスは EEVDF アルゴリズムでタスクを選ぶが，CFS の仮想実行時間の
計上はそのまま残っている．}
\begin{itemize}
  \item 常に最左の節点のタスク，すなわち仮想実行時間が最小のタスク，言い
    換えればこれまでに受け取った CPU 時間が最も少ないタスクをスケジュール
    する．
  \item タスクが実行されている間，CFS は定期的にその仮想実行時間を増やし，
    木の中の最左のタスクの仮想実行時間と比較する．
  \item 実行中のタスクの仮想実行時間の方がまだ小さければそのタスクは実行を
    続け，大きければそのタスクはプリエンプトされて，その仮想実行時間が
    定める位置に木へ挿入され，最左のタスクが実行される．
\end{itemize}

優先度は，仮想実行時間が進む速さによって表される．優先度の低いタスクでは
仮想実行時間が実時間より速く進むので，それらはすぐに CPU を失う．優先度の
高いタスクではより遅く進むので，それらはより長く実行される．こうして1本の
木がすべての優先度のタスクを保持する．スケジューラは最左の節点への
ポインタを保持しているのでタスクの選択は定数時間で済み，タスクの追加は
$N$ 個のタスクに対して $O(\log N)$ 時間で済む．

\begin{example}
  CFS はこの速さを nice 値ごとの重みから導く．nice 値 0 の重みは 1024 で
  あり，nice 値の1段階ごとに重みは約 1.25 倍変化する．nice 値 5 の重みは
  335 である．\tip{この係数は，nice 値1段が CPU のおよそ
  \SI{10}{\percent} に相当するように選ばれている．1段差なら
  $1.25/2.25 \approx 55:45$，5段差ならここでの 75:25 である．}重み $w$ のタスクが時間 $\Delta t$ だけ実行されると，その
  仮想実行時間は $\Delta t \cdot 1024/w$ だけ進む．タスク A（nice 値 0）
  とタスク B（nice 値 5）がそれぞれ \SI{6}{\milli\second} 実行されると，
  A の仮想実行時間は \SI{6}{\milli\second} 進むが，B のそれは
  $6 \cdot 1024/335 \approx \SI{18.3}{\milli\second}$ 進む．CFS は各
  タスクの仮想実行時間を近い値に保つので，B が実行される時間は A の約3分の
  1 となり，2つのタスクは CPU の $1024/1359 \approx \SI{75}{\percent}$ と
  \SI{25}{\percent} をそれぞれ受け取る．
\end{example}

\begin{table}[htbp]
  \centering
  \caption{Linux の2つのスケジューラ．}
  \label{tab:l07-linux}
  \small
  \begin{tabular}{@{}>{\raggedright\arraybackslash}p{0.22\linewidth}>{\raggedright\arraybackslash}p{0.34\linewidth}>{\raggedright\arraybackslash}p{0.34\linewidth}@{}}
    \toprule
    & $O(1)$ スケジューラ & CFS \\
    \midrule
    Linux のバージョン & 2.6.0--2.6.22 & 2.6.23--6.5 \\
    ランキュー & 140 本のリストからなるアクティブ配列と期限切れ配列 &
      仮想実行時間で順序付けた赤黒木 \\
    タスクの選択 & $O(1)$ & $O(1)$ \\
    タスクの追加 & $O(1)$ & $O(\log N)$ \\
    優先度 & 選択の順序，すなわち空でない最上位のリストと，タイムスライス
      の長さ．スリープ時間によるフィードバックが優先度を調整する &
      仮想実行時間が進む速さ \\
    対話的なタスク & アクティブ配列全体を待つことがある & 仮想実行時間が
      最小になり次第実行される \\
    \bottomrule
  \end{tabular}
\end{table}

\needspace{7\baselineskip}
\section{マルチプロセッサ上のスケジューリング}
\label{sec:l07-multi}

CPU が複数あると，スケジューリングはより複雑になる．\source{P3L1，
マルチプロセッサ上のスケジューリング；OSTEP ch.~10；Silberschatz ら
ch.~5．}2つのアーキテクチャを区別しなければならない
（\cref{fig:l07-multicpu}）．

\begin{definition}[共有メモリ型マルチプロセッサ]
  \term[きょうゆうめもりがたまるちぷろせっさ]{共有メモリ型マルチ%
  プロセッサ}[shared memory multiprocessor]は複数の CPU を持ち，各 CPU は
  自分専用のキャッシュを持ち，それらすべてが主記憶を共有する．
\end{definition}

\begin{definition}[マルチコアプロセッサ]
  \term[まるちこあぷろせっさ]{マルチコアプロセッサ}[multicore processor]は
  複数のコアを含み，各コアは自分専用のキャッシュを持つ CPU である．
  コアはプロセッサのラストレベルキャッシュ（LLC）と主記憶を共有する．
\end{definition}

\begin{figure}[htbp]
  \centering
  % (a) Shared memory multiprocessor: CPUs with private caches share memory.
% (b) Multicore processor: cores with private caches share the LLC.
% Usage: % (a) Shared memory multiprocessor: CPUs with private caches share memory.
% (b) Multicore processor: cores with private caches share the LLC.
% Usage: % (a) Shared memory multiprocessor: CPUs with private caches share memory.
% (b) Multicore processor: cores with private caches share the LLC.
% Usage: \input{figures/l07-multicpu}   (width ~112 mm)
\begin{tikzpicture}[x=1mm, y=1mm, font=\footnotesize\sffamily,
  cpu/.style={draw, fill=ln-accent!22, minimum width=17mm, minimum height=6mm, inner sep=1pt},
  cache/.style={draw, fill=ln-box, minimum width=17mm, minimum height=4.5mm,
                inner sep=1pt, font=\scriptsize\sffamily},
  mem/.style={draw, fill=ln-box, minimum height=7mm, inner sep=1pt},
  lab/.style={font=\scriptsize\sffamily, align=center}]
  % ---------------- (a) SMP
  \foreach \x/\n in {12/0, 40/1} {
    \node[cpu] (a\n) at (\x,0) {CPU\,\n};
    \node[cache, anchor=north] (al\n) at (a\n.south) {L1, L2};
    \node[cache, anchor=north] (ac\n) at (al\n.south) {LLC};
    \draw (ac\n.south) -- (\x,-20);
  }
  \draw[line width=1.2pt] (4,-20) -- (48,-20);
  \node[mem, minimum width=36mm] (am) at (26,-28) {\iflangja{メモリ（DRAM）}{memory (DRAM)}};
  \draw (am.north) -- (26,-20);
  \node[lab] at (26,-38) {\iflangja{(a) 共有メモリ型マルチプロセッサ}{(a) shared memory multiprocessor}};
  % ---------------- (b) multicore
  \draw[rounded corners=2pt, dashed, ln-muted] (64,4.5) rectangle (112,-17);
  \node[lab, anchor=north east, text=ln-muted] at (111.5,-17.5) {\iflangja{プロセッサ}{processor}};
  \foreach \x/\n in {76/0, 100/1} {
    \node[cpu] (b\n) at (\x,0) {\iflangja{コア}{core}\,\n};
    \node[cache, anchor=north] (bl\n) at (b\n.south) {L1, L2};
  }
  \node[cache, minimum width=41mm] (llc) at (88,-13) {\iflangja{LLC（共有）}{LLC (shared)}};
  \draw (bl0.south) -- (bl0.south |- llc.north);
  \draw (bl1.south) -- (bl1.south |- llc.north);
  \node[mem, minimum width=36mm] (bm) at (88,-28) {\iflangja{メモリ（DRAM）}{memory (DRAM)}};
  \draw (llc.south) -- (bm.north);
  \node[lab] at (88,-38) {\iflangja{(b) マルチコアプロセッサ}{(b) multicore processor}};
\end{tikzpicture}
   (width ~112 mm)
\begin{tikzpicture}[x=1mm, y=1mm, font=\footnotesize\sffamily,
  cpu/.style={draw, fill=ln-accent!22, minimum width=17mm, minimum height=6mm, inner sep=1pt},
  cache/.style={draw, fill=ln-box, minimum width=17mm, minimum height=4.5mm,
                inner sep=1pt, font=\scriptsize\sffamily},
  mem/.style={draw, fill=ln-box, minimum height=7mm, inner sep=1pt},
  lab/.style={font=\scriptsize\sffamily, align=center}]
  % ---------------- (a) SMP
  \foreach \x/\n in {12/0, 40/1} {
    \node[cpu] (a\n) at (\x,0) {CPU\,\n};
    \node[cache, anchor=north] (al\n) at (a\n.south) {L1, L2};
    \node[cache, anchor=north] (ac\n) at (al\n.south) {LLC};
    \draw (ac\n.south) -- (\x,-20);
  }
  \draw[line width=1.2pt] (4,-20) -- (48,-20);
  \node[mem, minimum width=36mm] (am) at (26,-28) {\iflangja{メモリ（DRAM）}{memory (DRAM)}};
  \draw (am.north) -- (26,-20);
  \node[lab] at (26,-38) {\iflangja{(a) 共有メモリ型マルチプロセッサ}{(a) shared memory multiprocessor}};
  % ---------------- (b) multicore
  \draw[rounded corners=2pt, dashed, ln-muted] (64,4.5) rectangle (112,-17);
  \node[lab, anchor=north east, text=ln-muted] at (111.5,-17.5) {\iflangja{プロセッサ}{processor}};
  \foreach \x/\n in {76/0, 100/1} {
    \node[cpu] (b\n) at (\x,0) {\iflangja{コア}{core}\,\n};
    \node[cache, anchor=north] (bl\n) at (b\n.south) {L1, L2};
  }
  \node[cache, minimum width=41mm] (llc) at (88,-13) {\iflangja{LLC（共有）}{LLC (shared)}};
  \draw (bl0.south) -- (bl0.south |- llc.north);
  \draw (bl1.south) -- (bl1.south |- llc.north);
  \node[mem, minimum width=36mm] (bm) at (88,-28) {\iflangja{メモリ（DRAM）}{memory (DRAM)}};
  \draw (llc.south) -- (bm.north);
  \node[lab] at (88,-38) {\iflangja{(b) マルチコアプロセッサ}{(b) multicore processor}};
\end{tikzpicture}
   (width ~112 mm)
\begin{tikzpicture}[x=1mm, y=1mm, font=\footnotesize\sffamily,
  cpu/.style={draw, fill=ln-accent!22, minimum width=17mm, minimum height=6mm, inner sep=1pt},
  cache/.style={draw, fill=ln-box, minimum width=17mm, minimum height=4.5mm,
                inner sep=1pt, font=\scriptsize\sffamily},
  mem/.style={draw, fill=ln-box, minimum height=7mm, inner sep=1pt},
  lab/.style={font=\scriptsize\sffamily, align=center}]
  % ---------------- (a) SMP
  \foreach \x/\n in {12/0, 40/1} {
    \node[cpu] (a\n) at (\x,0) {CPU\,\n};
    \node[cache, anchor=north] (al\n) at (a\n.south) {L1, L2};
    \node[cache, anchor=north] (ac\n) at (al\n.south) {LLC};
    \draw (ac\n.south) -- (\x,-20);
  }
  \draw[line width=1.2pt] (4,-20) -- (48,-20);
  \node[mem, minimum width=36mm] (am) at (26,-28) {\iflangja{メモリ（DRAM）}{memory (DRAM)}};
  \draw (am.north) -- (26,-20);
  \node[lab] at (26,-38) {\iflangja{(a) 共有メモリ型マルチプロセッサ}{(a) shared memory multiprocessor}};
  % ---------------- (b) multicore
  \draw[rounded corners=2pt, dashed, ln-muted] (64,4.5) rectangle (112,-17);
  \node[lab, anchor=north east, text=ln-muted] at (111.5,-17.5) {\iflangja{プロセッサ}{processor}};
  \foreach \x/\n in {76/0, 100/1} {
    \node[cpu] (b\n) at (\x,0) {\iflangja{コア}{core}\,\n};
    \node[cache, anchor=north] (bl\n) at (b\n.south) {L1, L2};
  }
  \node[cache, minimum width=41mm] (llc) at (88,-13) {\iflangja{LLC（共有）}{LLC (shared)}};
  \draw (bl0.south) -- (bl0.south |- llc.north);
  \draw (bl1.south) -- (bl1.south |- llc.north);
  \node[mem, minimum width=36mm] (bm) at (88,-28) {\iflangja{メモリ（DRAM）}{memory (DRAM)}};
  \draw (llc.south) -- (bm.north);
  \node[lab] at (88,-38) {\iflangja{(b) マルチコアプロセッサ}{(b) multicore processor}};
\end{tikzpicture}

  \caption{(a) 共有メモリ型マルチプロセッサ．どの CPU も自分専用の
    キャッシュを持つ．(b) マルチコアプロセッサ．コアは専用のキャッシュを
    持ち，ラストレベルキャッシュを共有する．}
  \label{fig:l07-multicpu}
\end{figure}

オペレーティングシステムから見れば，マルチコアプロセッサの各コアは1つの
CPU であり，システムは複数のマルチコアプロセッサを持つことで両方の
アーキテクチャを組み合わせることもある．

\subsection{キャッシュアフィニティと負荷分散}

タスクの性能はキャッシュの状態に依存する．タスクがある CPU 上で実行される
間，その CPU のキャッシュはそのタスクのデータで満たされる．次回も同じ CPU
上で実行されればホットキャッシュを見出すが（第2章），別の CPU 上で実行
されればコールドキャッシュから始まり，そのキャッシュが満たされるまで実行
は遅くなる．\margin{直近 \texttt{sched\_migration\_cost\_ns}（既定は
\SI{500}{\micro\second}）以内に実行されたタスクはキャッシュがホットと
みなされ，負荷分散はそれを動かさない．}

\begin{definition}[キャッシュアフィニティ]
  タスクが，自分のデータを保持するキャッシュを持つ CPU 上で実行されること
  から得る利益を，その\term[きゃっしゅあふぃにてぃ]{キャッシュ%
  アフィニティ}[cache affinity]という．スケジューラは，各タスクをできる
  限り同じ CPU 上に保つことでこれを維持する．
\end{definition}

そこでマルチプロセッサのためのスケジューラは階層的なアーキテクチャを取る．
各 CPU が自分のランキューを持ち，その CPU のためのスケジューラはそこから
のみタスクを選ぶので，タスクは同じ CPU 上で実行され続ける．CPU ごとの
スケジューラの上には\term[ふかぶんさん]{負荷分散}[load balancing]の構成
要素があり，ある CPU が過負荷である一方で別の CPU の仕事が少なすぎる，と
いうことがないようにランキューの間でタスクを移動させる．負荷分散は
ランキューの長さに基づいて動作するほか，ある CPU がアイドルになったときに
他の CPU のランキューからタスクを取ってくることもある．\margin{Linux では，
アプリケーションは \texttt{sched\_setaffinity} でスレッドを CPU の集合に
限定できる．}

\subsection{NUMA 対応スケジューリング}

マルチプロセッサは複数のメモリノードを持つこともある
（\cref{fig:l07-numa}）．

\begin{definition}[NUMA]
  \term{NUMA}[non-uniform memory access]%
  \index{non-uniform memory access|see{NUMA}}（不均一メモリアクセス）を
  持つシステムでは，主記憶は複数のメモリノードからなり，各ノードは1つ
  以上のプロセッサのソケットの近くに置かれる．CPU から自分のローカルな
  メモリノードへのアクセスは，ソケット間のインターコネクト，例えば Intel
  の QuickPath Interconnect を経由するリモートのノードへのアクセスより
  速い．
\end{definition}

\begin{figure}[htbp]
  \centering
  % Hierarchical scheduling on a NUMA system: per-CPU runqueues, a load
% balancer, and two memory nodes connected by an interconnect.
% Usage: % Hierarchical scheduling on a NUMA system: per-CPU runqueues, a load
% balancer, and two memory nodes connected by an interconnect.
% Usage: % Hierarchical scheduling on a NUMA system: per-CPU runqueues, a load
% balancer, and two memory nodes connected by an interconnect.
% Usage: \input{figures/l07-numa}   (width ~112 mm)
\begin{tikzpicture}[x=1mm, y=1mm, font=\footnotesize\sffamily,
  >={Stealth[length=1.6mm]},
  cpu/.style={draw, fill=ln-accent!22, minimum width=16mm, minimum height=6mm, inner sep=1pt},
  mem/.style={draw, fill=ln-box, minimum width=38mm, minimum height=7mm, inner sep=1pt},
  nd/.style={draw, dashed, rounded corners=2pt, ln-muted},
  lab/.style={font=\scriptsize\sffamily, align=center},
  tk/.style={draw, circle, fill=white, minimum size=2.6mm, inner sep=0pt}]
  % load balancer
  \draw[fill=ln-box] (0,36) rectangle (112,42);
  \node at (56,39) {\iflangja{負荷分散}{load balancing}};
  % per-CPU runqueues and CPUs
  \foreach \x/\n/\k in {13/0/3, 36/1/1, 76/2/2, 99/3/0} {
    \draw[<->] (\x,35.5) -- (\x,29.5);
    \draw[fill=white] (\x-8,24) rectangle (\x+8,29);
    \foreach \i in {1,...,3} \draw (\x+8-4*\i,24) -- ++(0,5);
    \ifnum\k>0 \foreach \i in {1,...,\k} \node[tk] at (\x+10-4*\i,26.5) {}; \fi
    \draw[->] (\x,24) -- (\x,17);
    \node[cpu] (c\n) at (\x,14) {CPU\,\n};
  }
  \node[lab] at (56,26.5) {\iflangja{CPU ごとの\\ランキュー}{per-CPU\\runqueues}};
  % memory nodes
  \node[mem] (m0) at (24.5,-6) {\iflangja{メモリノード 0}{memory node 0}};
  \node[mem] (m1) at (87.5,-6) {\iflangja{メモリノード 1}{memory node 1}};
  \draw[nd] (2,20) rectangle (47,-12);
  \draw[nd] (65,20) rectangle (110,-12);
  \node[lab, anchor=north, text=ln-muted] at (24.5,-12.5) {\iflangja{ノード 0（ソケット 0）}{node 0 (socket 0)}};
  \node[lab, anchor=north, text=ln-muted] at (87.5,-12.5) {\iflangja{ノード 1（ソケット 1）}{node 1 (socket 1)}};
  % local access
  \draw[<->, thick] (c0.south) -- node[lab, right] {\iflangja{ローカル：\\速い}{local:\\fast}} (c0.south |- m0.north);
  \draw (c1.south) -- (c1.south |- m0.north);
  \draw (c2.south) -- (c2.south |- m1.north);
  \draw (c3.south) -- (c3.south |- m1.north);
  % interconnect
  \draw[line width=2pt, ln-muted] (m0.east) -- (m1.west);
  \node[lab, anchor=north, text=ln-muted] at (56,-12.5) {\iflangja{インターコネクト}{interconnect}};
  \draw[ln-muted, densely dotted] (56,-6) -- (56,-12.3);
  % remote access
  \draw[<->, thick, densely dashed, ln-caution] (c2.south west) -- (m0.north east);
  \node[lab, text=ln-caution] at (53.5,10.5) {\iflangja{リモート：\\遅い}{remote:\\slow}};
\end{tikzpicture}
   (width ~112 mm)
\begin{tikzpicture}[x=1mm, y=1mm, font=\footnotesize\sffamily,
  >={Stealth[length=1.6mm]},
  cpu/.style={draw, fill=ln-accent!22, minimum width=16mm, minimum height=6mm, inner sep=1pt},
  mem/.style={draw, fill=ln-box, minimum width=38mm, minimum height=7mm, inner sep=1pt},
  nd/.style={draw, dashed, rounded corners=2pt, ln-muted},
  lab/.style={font=\scriptsize\sffamily, align=center},
  tk/.style={draw, circle, fill=white, minimum size=2.6mm, inner sep=0pt}]
  % load balancer
  \draw[fill=ln-box] (0,36) rectangle (112,42);
  \node at (56,39) {\iflangja{負荷分散}{load balancing}};
  % per-CPU runqueues and CPUs
  \foreach \x/\n/\k in {13/0/3, 36/1/1, 76/2/2, 99/3/0} {
    \draw[<->] (\x,35.5) -- (\x,29.5);
    \draw[fill=white] (\x-8,24) rectangle (\x+8,29);
    \foreach \i in {1,...,3} \draw (\x+8-4*\i,24) -- ++(0,5);
    \ifnum\k>0 \foreach \i in {1,...,\k} \node[tk] at (\x+10-4*\i,26.5) {}; \fi
    \draw[->] (\x,24) -- (\x,17);
    \node[cpu] (c\n) at (\x,14) {CPU\,\n};
  }
  \node[lab] at (56,26.5) {\iflangja{CPU ごとの\\ランキュー}{per-CPU\\runqueues}};
  % memory nodes
  \node[mem] (m0) at (24.5,-6) {\iflangja{メモリノード 0}{memory node 0}};
  \node[mem] (m1) at (87.5,-6) {\iflangja{メモリノード 1}{memory node 1}};
  \draw[nd] (2,20) rectangle (47,-12);
  \draw[nd] (65,20) rectangle (110,-12);
  \node[lab, anchor=north, text=ln-muted] at (24.5,-12.5) {\iflangja{ノード 0（ソケット 0）}{node 0 (socket 0)}};
  \node[lab, anchor=north, text=ln-muted] at (87.5,-12.5) {\iflangja{ノード 1（ソケット 1）}{node 1 (socket 1)}};
  % local access
  \draw[<->, thick] (c0.south) -- node[lab, right] {\iflangja{ローカル：\\速い}{local:\\fast}} (c0.south |- m0.north);
  \draw (c1.south) -- (c1.south |- m0.north);
  \draw (c2.south) -- (c2.south |- m1.north);
  \draw (c3.south) -- (c3.south |- m1.north);
  % interconnect
  \draw[line width=2pt, ln-muted] (m0.east) -- (m1.west);
  \node[lab, anchor=north, text=ln-muted] at (56,-12.5) {\iflangja{インターコネクト}{interconnect}};
  \draw[ln-muted, densely dotted] (56,-6) -- (56,-12.3);
  % remote access
  \draw[<->, thick, densely dashed, ln-caution] (c2.south west) -- (m0.north east);
  \node[lab, text=ln-caution] at (53.5,10.5) {\iflangja{リモート：\\遅い}{remote:\\slow}};
\end{tikzpicture}
   (width ~112 mm)
\begin{tikzpicture}[x=1mm, y=1mm, font=\footnotesize\sffamily,
  >={Stealth[length=1.6mm]},
  cpu/.style={draw, fill=ln-accent!22, minimum width=16mm, minimum height=6mm, inner sep=1pt},
  mem/.style={draw, fill=ln-box, minimum width=38mm, minimum height=7mm, inner sep=1pt},
  nd/.style={draw, dashed, rounded corners=2pt, ln-muted},
  lab/.style={font=\scriptsize\sffamily, align=center},
  tk/.style={draw, circle, fill=white, minimum size=2.6mm, inner sep=0pt}]
  % load balancer
  \draw[fill=ln-box] (0,36) rectangle (112,42);
  \node at (56,39) {\iflangja{負荷分散}{load balancing}};
  % per-CPU runqueues and CPUs
  \foreach \x/\n/\k in {13/0/3, 36/1/1, 76/2/2, 99/3/0} {
    \draw[<->] (\x,35.5) -- (\x,29.5);
    \draw[fill=white] (\x-8,24) rectangle (\x+8,29);
    \foreach \i in {1,...,3} \draw (\x+8-4*\i,24) -- ++(0,5);
    \ifnum\k>0 \foreach \i in {1,...,\k} \node[tk] at (\x+10-4*\i,26.5) {}; \fi
    \draw[->] (\x,24) -- (\x,17);
    \node[cpu] (c\n) at (\x,14) {CPU\,\n};
  }
  \node[lab] at (56,26.5) {\iflangja{CPU ごとの\\ランキュー}{per-CPU\\runqueues}};
  % memory nodes
  \node[mem] (m0) at (24.5,-6) {\iflangja{メモリノード 0}{memory node 0}};
  \node[mem] (m1) at (87.5,-6) {\iflangja{メモリノード 1}{memory node 1}};
  \draw[nd] (2,20) rectangle (47,-12);
  \draw[nd] (65,20) rectangle (110,-12);
  \node[lab, anchor=north, text=ln-muted] at (24.5,-12.5) {\iflangja{ノード 0（ソケット 0）}{node 0 (socket 0)}};
  \node[lab, anchor=north, text=ln-muted] at (87.5,-12.5) {\iflangja{ノード 1（ソケット 1）}{node 1 (socket 1)}};
  % local access
  \draw[<->, thick] (c0.south) -- node[lab, right] {\iflangja{ローカル：\\速い}{local:\\fast}} (c0.south |- m0.north);
  \draw (c1.south) -- (c1.south |- m0.north);
  \draw (c2.south) -- (c2.south |- m1.north);
  \draw (c3.south) -- (c3.south |- m1.north);
  % interconnect
  \draw[line width=2pt, ln-muted] (m0.east) -- (m1.west);
  \node[lab, anchor=north, text=ln-muted] at (56,-12.5) {\iflangja{インターコネクト}{interconnect}};
  \draw[ln-muted, densely dotted] (56,-6) -- (56,-12.3);
  % remote access
  \draw[<->, thick, densely dashed, ln-caution] (c2.south west) -- (m0.north east);
  \node[lab, text=ln-caution] at (53.5,10.5) {\iflangja{リモート：\\遅い}{remote:\\slow}};
\end{tikzpicture}

  \caption{NUMA システム上の階層的なスケジューリング．どの CPU も自分の
    ランキューからタスクをスケジュールし，負荷分散がランキューの間で
    タスクを移動させる．CPU は自分のノードのメモリには直接，もう一方の
    ノードのメモリにはインターコネクトを経由して到達する．}
  \label{fig:l07-numa}
\end{figure}

\term[NUMAたいおうすけじゅーりんぐ]{NUMA対応スケジューリング}%
[NUMA-aware scheduling]では，スケジューラはタスクを，その状態を保持する
メモリノードに近い CPU 上に保ち，メモリアクセスがローカルなままであるよう
にする．負荷分散はこれを考慮しなければならない．タスクを別のノードの CPU
へ移すとランキューの均衡は取れるかもしれないが，そのタスクの主記憶への
アクセスはすべてインターコネクトを越えることになり，遅くなる．

\begin{example}
  あるタスクが実行時間の \SI{40}{\percent} を主記憶へのアクセスに費やし，
  リモートのアクセスはローカルのアクセスの 1.5 倍の時間を要するとする．
  負荷分散器がこのタスクを，そのメモリはそのままにして，もう一方の
  ノードの CPU へ移すと，そのタスクの実行時間は
  $0.6 + 0.4 \times 1.5 = 1.2$ 倍，すなわち \SI{20}{\percent} 増える．
\end{example}

\needspace{7\baselineskip}
\section{ハイパースレッディング}
\label{sec:l07-smt}

CPU を複数のタスクで多重化しなければならないのは，タスクの実行コンテキスト，
すなわちプログラムカウンタやスタックポインタといったそのレジスタの値が，
タスクの実行中は CPU のレジスタに保持されるからである．\source{P3L1，
ハイパースレッディング；Silberschatz ら ch.~5；Fedorova ら
\cite{fedorova2004}．}別のタスクに切り替えるには，そのコンテキストを退避
し，別のコンテキストを読み込まなければならない．その代わりに，ハードウェア
が1つの CPU 上に複数のレジスタ集合，したがって複数のハードウェアで支援
された実行コンテキストを提供することもできる．その CPU は，コンテキストに
共有される1組の実行資源を持つ単一の CPU のままであるが，あるコンテキスト
から別のコンテキストへ移るのに状態の退避と復元を必要としない．

\begin{definition}[同時マルチスレッディング]
  \term[どうじまるちすれっでぃんぐ]{同時マルチスレッディング}%
  [simultaneous multithreading]%
  \index{SMT|see{同時マルチスレッディング}}%
  \index{はいぱーすれっでぃんぐ@ハイパースレッディング|see{同時マルチスレッディング}}（SMT）
  では，1つの CPU が複数のハードウェア実行コンテキスト，すなわち
  ハードウェアスレッドを提供する．それぞれが自分のレジスタ状態を持ち，
  CPU のパイプラインとキャッシュを共有する．
\end{definition}

この技法は，ハードウェアマルチスレッディング，Intel が用いる名称である
ハイパースレッディング，あるいはチップマルチスレッディングとも呼ばれる．
より正確には，\term[ちっぷまるちすれっでぃんぐ]{チップマルチスレッディング}%
[chip multithreading]\index{CMT|see{チップマルチスレッディング}}（CMT）は
複数のコアとコア当たり複数のハードウェアスレッドを組み合わせたプロセッサ
を指し，これは Fedorova らが研究したプラットフォームである．実際にはこれら
の名称はしばしば同じ意味で使われる．コア当たり8つまでのハードウェアスレッド
を持つプロセッサも作られてきたが，最も一般的な構成は2つであり，SMT は起動
時に有効にも無効にもできる．\margin{Intel の最初の実装は，ダイ面積の
\SI{5}{\percent} 未満の増加でスループットをおよそ \SI{25}{\percent}
高めた（Marr ら，\emph{Intel Technology Journal}，2002）．}

設計によって，コンテキストがパイプラインをどれだけ密に共有するかは異なる．
最も単純なものはコンテキストを交互に実行し，短ければ1サイクルごとに
パイプラインを別のコンテキストへ渡す．狭義の同時マルチスレッディング，
例えば Intel の Hyper-Threading では，複数のハードウェアスレッドの命令が同じ
サイクルにパイプラインの各段を占める．\margin{Fedorova らがシミュレート
したコアは，4つのハードウェアスレッドを1サイクルごとに交互に実行する．}

各ハードウェアスレッドは，オペレーティングシステムには個別の仮想 CPU と
して見え，スケジューラはその上にタスクを配置できる．したがって
スケジューラは，同じ CPU のハードウェアスレッド上でどのタスクを一緒に
スケジュールするかを決めなければならない．ハードウェアスレッド間の
コンテキストスイッチは数サイクル程度であるのに対し，キャッシュミスの後に
メモリからデータを読み込むには 100 サイクル程度を要する．第3章の議論に
よれば，切り替えが有利になるのはアイドル時間が切り替えの費用の2倍を超える
ときであり（\cref{eq:l03-idle}），ここではそれが大きな余裕をもって成り立つ．
したがって SMT により，1つの CPU 上のマルチスレッドが I/O のレイテンシを
隠すのとちょうど同じように，CPU はメモリアクセスのレイテンシを隠せる．

\needspace{7\baselineskip}
\section{ハイパースレッディング対応プラットフォームのスケジューリング}
\label{sec:l07-smt-sched}

どのタスクが CPU を共有すべきだろうか．Fedorova らが研究した
プラットフォームに倣ったマルチスレッド CPU の単純なモデルが，この問いに
答える．\source{P3L1，ハイパースレッディング対応プラットフォームの
スケジューリング；Fedorova ら \cite{fedorova2004}．}このモデルは4つの
仮定に立つ．
\begin{enumerate}
  \item スレッドは CPU の毎サイクルに1つの命令を発行し，パイプラインは
    1サイクル当たり高々1つの命令を実行する．
  \item メモリバウンドのスレッドのメモリアクセスは4サイクルを要し，その
    間そのスレッドは命令を発行しない．この小さな値はタイミング図を短く
    保つためのものであり，メモリまで到達するアクセスの 100 サイクル程度
    ではなく，L1 キャッシュで処理されるロードの遅延に当たる．
  \item ハードウェアスレッド間の切り替えは瞬時に行われる．
  \item CPU は2つのハードウェアスレッドを持つ．
\end{enumerate}
このモデルでは，CPU バウンドのタスクは毎サイクルにパイプラインを必要と
する．\term[めもりばうんどたすく]{メモリバウンドタスク}[memory-bound task]
とは，その進行がメモリアクセスによって制限されるタスクであり，ほとんどの
サイクルをメモリからのデータの到着を待って費やす．\cref{fig:l07-smt}は，
ありうる3通りの組み合わせを示している．

\begin{figure}[htbp]
  \centering
  % Co-scheduling two threads on the two hardware threads of an SMT CPU.
% One instruction per cycle on the shared pipeline; a memory access takes
% 4 cycles; switching between hardware threads is instantaneous.
% Usage: % Co-scheduling two threads on the two hardware threads of an SMT CPU.
% One instruction per cycle on the shared pipeline; a memory access takes
% 4 cycles; switching between hardware threads is instantaneous.
% Usage: % Co-scheduling two threads on the two hardware threads of an SMT CPU.
% One instruction per cycle on the shared pipeline; a memory access takes
% 4 cycles; switching between hardware threads is instantaneous.
% Usage: \input{figures/l07-smt}   (width ~104 mm)
\begin{tikzpicture}[x=8mm, y=1mm, font=\scriptsize\sffamily,
  ex/.style={draw, fill=ln-accent!30},
  mw/.style={draw, fill=ln-box},
  rd/.style={draw, fill=white},
  lab/.style={anchor=east, font=\scriptsize\sffamily},
  ttl/.style={anchor=south west, font=\footnotesize\sffamily\bfseries, inner sep=1pt}]
  % cycle numbers
  \node[lab] at (-0.1,60) {\iflangja{サイクル}{cycle}};
  \foreach \c in {1,...,10} \node at (\c-0.5,60) {\c};
  % cells: \cell{row y}{cycle}{style}{text}
  \def\cell#1#2#3#4{\draw[#3] (#2-1,#1) rectangle (#2,#1+5); \node at (#2-0.5,#1+2.5) {#4};}
  % ---------- (a) two CPU-bound threads
  \node[ttl] at (0,51.5) {\iflangja{(a) CPU バウンド 2 つ}{(a) two CPU-bound threads}};
  \node[lab] at (-0.1,48.5) {HT\,0: C};
  \node[lab] at (-0.1,42.5) {HT\,1: C};
  \foreach \c in {1,3,5,7,9} { \cell{46}{\c}{ex}{C} \cell{40}{\c}{rd}{} }
  \foreach \c in {2,4,6,8,10} { \cell{46}{\c}{rd}{} \cell{40}{\c}{ex}{C} }
  % ---------- (b) two memory-bound threads
  \node[ttl] at (0,32.5) {\iflangja{(b) メモリバウンド 2 つ}{(b) two memory-bound threads}};
  \node[lab] at (-0.1,29.5) {HT\,0: M};
  \node[lab] at (-0.1,23.5) {HT\,1: M};
  \foreach \c in {1,6} { \cell{27}{\c}{ex}{M} }
  \foreach \c in {2,3,4,5,7,8,9,10} { \cell{27}{\c}{mw}{} }
  \foreach \c in {2,7} { \cell{21}{\c}{ex}{M} }
  \foreach \c in {1,3,4,5,6,8,9,10} { \cell{21}{\c}{mw}{} }
  \foreach \c in {3,4,5,8,9,10} { \node[text=ln-caution] at (\c-0.5,18.6) {\iflangja{空き}{idle}}; }
  % ---------- (c) one CPU-bound and one memory-bound thread
  \node[ttl] at (0,9.5) {\iflangja{(c) CPU バウンドとメモリバウンド}{(c) one CPU-bound, one memory-bound thread}};
  \node[lab] at (-0.1,6.5) {HT\,0: M};
  \node[lab] at (-0.1,0.5) {HT\,1: C};
  \foreach \c in {1,6} { \cell{4}{\c}{ex}{M} }
  \foreach \c in {2,3,4,5,7,8,9,10} { \cell{4}{\c}{mw}{} }
  \foreach \c in {1,6} { \cell{-2}{\c}{rd}{} }
  \foreach \c in {2,3,4,5,7,8,9,10} { \cell{-2}{\c}{ex}{C} }
  % ---------- legend
  \begin{scope}[yshift=-9mm]
    \draw[ex] (0,-1.5) rectangle (0.5,1.5);
    \node[anchor=west] at (0.55,0) {\iflangja{命令を実行}{executes an instruction}};
    \draw[mw] (3.9,-1.5) rectangle (4.4,1.5);
    \node[anchor=west] at (4.45,0) {\iflangja{メモリを待つ}{waits for memory}};
    \draw[rd] (7.2,-1.5) rectangle (7.7,1.5);
    \node[anchor=west] at (7.75,0) {\iflangja{実行可能だが待つ}{ready, not executing}};
  \end{scope}
\end{tikzpicture}
   (width ~104 mm)
\begin{tikzpicture}[x=8mm, y=1mm, font=\scriptsize\sffamily,
  ex/.style={draw, fill=ln-accent!30},
  mw/.style={draw, fill=ln-box},
  rd/.style={draw, fill=white},
  lab/.style={anchor=east, font=\scriptsize\sffamily},
  ttl/.style={anchor=south west, font=\footnotesize\sffamily\bfseries, inner sep=1pt}]
  % cycle numbers
  \node[lab] at (-0.1,60) {\iflangja{サイクル}{cycle}};
  \foreach \c in {1,...,10} \node at (\c-0.5,60) {\c};
  % cells: \cell{row y}{cycle}{style}{text}
  \def\cell#1#2#3#4{\draw[#3] (#2-1,#1) rectangle (#2,#1+5); \node at (#2-0.5,#1+2.5) {#4};}
  % ---------- (a) two CPU-bound threads
  \node[ttl] at (0,51.5) {\iflangja{(a) CPU バウンド 2 つ}{(a) two CPU-bound threads}};
  \node[lab] at (-0.1,48.5) {HT\,0: C};
  \node[lab] at (-0.1,42.5) {HT\,1: C};
  \foreach \c in {1,3,5,7,9} { \cell{46}{\c}{ex}{C} \cell{40}{\c}{rd}{} }
  \foreach \c in {2,4,6,8,10} { \cell{46}{\c}{rd}{} \cell{40}{\c}{ex}{C} }
  % ---------- (b) two memory-bound threads
  \node[ttl] at (0,32.5) {\iflangja{(b) メモリバウンド 2 つ}{(b) two memory-bound threads}};
  \node[lab] at (-0.1,29.5) {HT\,0: M};
  \node[lab] at (-0.1,23.5) {HT\,1: M};
  \foreach \c in {1,6} { \cell{27}{\c}{ex}{M} }
  \foreach \c in {2,3,4,5,7,8,9,10} { \cell{27}{\c}{mw}{} }
  \foreach \c in {2,7} { \cell{21}{\c}{ex}{M} }
  \foreach \c in {1,3,4,5,6,8,9,10} { \cell{21}{\c}{mw}{} }
  \foreach \c in {3,4,5,8,9,10} { \node[text=ln-caution] at (\c-0.5,18.6) {\iflangja{空き}{idle}}; }
  % ---------- (c) one CPU-bound and one memory-bound thread
  \node[ttl] at (0,9.5) {\iflangja{(c) CPU バウンドとメモリバウンド}{(c) one CPU-bound, one memory-bound thread}};
  \node[lab] at (-0.1,6.5) {HT\,0: M};
  \node[lab] at (-0.1,0.5) {HT\,1: C};
  \foreach \c in {1,6} { \cell{4}{\c}{ex}{M} }
  \foreach \c in {2,3,4,5,7,8,9,10} { \cell{4}{\c}{mw}{} }
  \foreach \c in {1,6} { \cell{-2}{\c}{rd}{} }
  \foreach \c in {2,3,4,5,7,8,9,10} { \cell{-2}{\c}{ex}{C} }
  % ---------- legend
  \begin{scope}[yshift=-9mm]
    \draw[ex] (0,-1.5) rectangle (0.5,1.5);
    \node[anchor=west] at (0.55,0) {\iflangja{命令を実行}{executes an instruction}};
    \draw[mw] (3.9,-1.5) rectangle (4.4,1.5);
    \node[anchor=west] at (4.45,0) {\iflangja{メモリを待つ}{waits for memory}};
    \draw[rd] (7.2,-1.5) rectangle (7.7,1.5);
    \node[anchor=west] at (7.75,0) {\iflangja{実行可能だが待つ}{ready, not executing}};
  \end{scope}
\end{tikzpicture}
   (width ~104 mm)
\begin{tikzpicture}[x=8mm, y=1mm, font=\scriptsize\sffamily,
  ex/.style={draw, fill=ln-accent!30},
  mw/.style={draw, fill=ln-box},
  rd/.style={draw, fill=white},
  lab/.style={anchor=east, font=\scriptsize\sffamily},
  ttl/.style={anchor=south west, font=\footnotesize\sffamily\bfseries, inner sep=1pt}]
  % cycle numbers
  \node[lab] at (-0.1,60) {\iflangja{サイクル}{cycle}};
  \foreach \c in {1,...,10} \node at (\c-0.5,60) {\c};
  % cells: \cell{row y}{cycle}{style}{text}
  \def\cell#1#2#3#4{\draw[#3] (#2-1,#1) rectangle (#2,#1+5); \node at (#2-0.5,#1+2.5) {#4};}
  % ---------- (a) two CPU-bound threads
  \node[ttl] at (0,51.5) {\iflangja{(a) CPU バウンド 2 つ}{(a) two CPU-bound threads}};
  \node[lab] at (-0.1,48.5) {HT\,0: C};
  \node[lab] at (-0.1,42.5) {HT\,1: C};
  \foreach \c in {1,3,5,7,9} { \cell{46}{\c}{ex}{C} \cell{40}{\c}{rd}{} }
  \foreach \c in {2,4,6,8,10} { \cell{46}{\c}{rd}{} \cell{40}{\c}{ex}{C} }
  % ---------- (b) two memory-bound threads
  \node[ttl] at (0,32.5) {\iflangja{(b) メモリバウンド 2 つ}{(b) two memory-bound threads}};
  \node[lab] at (-0.1,29.5) {HT\,0: M};
  \node[lab] at (-0.1,23.5) {HT\,1: M};
  \foreach \c in {1,6} { \cell{27}{\c}{ex}{M} }
  \foreach \c in {2,3,4,5,7,8,9,10} { \cell{27}{\c}{mw}{} }
  \foreach \c in {2,7} { \cell{21}{\c}{ex}{M} }
  \foreach \c in {1,3,4,5,6,8,9,10} { \cell{21}{\c}{mw}{} }
  \foreach \c in {3,4,5,8,9,10} { \node[text=ln-caution] at (\c-0.5,18.6) {\iflangja{空き}{idle}}; }
  % ---------- (c) one CPU-bound and one memory-bound thread
  \node[ttl] at (0,9.5) {\iflangja{(c) CPU バウンドとメモリバウンド}{(c) one CPU-bound, one memory-bound thread}};
  \node[lab] at (-0.1,6.5) {HT\,0: M};
  \node[lab] at (-0.1,0.5) {HT\,1: C};
  \foreach \c in {1,6} { \cell{4}{\c}{ex}{M} }
  \foreach \c in {2,3,4,5,7,8,9,10} { \cell{4}{\c}{mw}{} }
  \foreach \c in {1,6} { \cell{-2}{\c}{rd}{} }
  \foreach \c in {2,3,4,5,7,8,9,10} { \cell{-2}{\c}{ex}{C} }
  % ---------- legend
  \begin{scope}[yshift=-9mm]
    \draw[ex] (0,-1.5) rectangle (0.5,1.5);
    \node[anchor=west] at (0.55,0) {\iflangja{命令を実行}{executes an instruction}};
    \draw[mw] (3.9,-1.5) rectangle (4.4,1.5);
    \node[anchor=west] at (4.45,0) {\iflangja{メモリを待つ}{waits for memory}};
    \draw[rd] (7.2,-1.5) rectangle (7.7,1.5);
    \node[anchor=west] at (7.75,0) {\iflangja{実行可能だが待つ}{ready, not executing}};
  \end{scope}
\end{tikzpicture}

  \caption{1つの SMT CPU のハードウェアスレッド（HT）上に2つのスレッドを
    一緒にスケジュールした場合．C は CPU バウンドの，M はメモリバウンドの
    スレッドを表す．各列は共有パイプラインの1サイクルである．}
  \label{fig:l07-smt}
\end{figure}

\begin{description}
  \item[CPU バウンドのスレッド2つ．] どちらのスレッドも毎サイクル実行
    可能だが，パイプラインは1サイクル当たり1命令しか実行しない．2つの
    スレッドは干渉し合ってパイプラインを奪い合う．それぞれは単独のときの
    半分の速度でしか進まず，2倍の性能低下となる．その間，メモリの側は
    アイドルである．
  \item[メモリバウンドのスレッド2つ．] どちらのスレッドもほとんどの
    サイクルをメモリ待ちで過ごし，図の 10 サイクルのうち 6 サイクルでは
    どちらも実行すべき命令を持たない．これらの CPU サイクルは無駄になる．
  \item[CPU バウンドとメモリバウンドのスレッド．] CPU バウンドのスレッド
    は，メモリバウンドのスレッドが待っているサイクルを使う．すなわち，
    10 サイクルのうち 8 サイクルで実行し，その間メモリバウンドのスレッドは
    単独のときと同じ速さで進む．アイドルのサイクルはなく，CPU もメモリも
    よく活用される．CPU バウンドのスレッドは5サイクルに1つを失うので干渉
    は残るが，それは小さい．
\end{description}

\needspace{6\baselineskip}
\begin{keypoint}
  SMT の CPU では，CPU バウンドのタスクとメモリバウンドのタスクを一緒に
  スケジュールすることで，パイプラインの奪い合いとアイドルサイクルの
  両方を避けられる．したがってスケジューラは，各 CPU のハードウェア
  スレッド上で2種類のタスクを混ぜるべきである．
\end{keypoint}

\needspace{7\baselineskip}
\section{ハードウェア性能カウンタ}
\label{sec:l07-counters}

CPU バウンドのタスクとメモリバウンドのタスクを混ぜるには，スケジューラが
どちらがどちらであるかを知る必要がある．\source{P3L1，CPU バウンドか
メモリバウンドか，ハードウェアカウンタを用いたスケジューリング；
Fedorova ら．}I/O バウンドのタスクに用いた履歴に基づく経験則は通用しない．
$O(1)$ スケジューラが用いるスリープ時間は，メモリバウンドのタスクを明らか
にしない．メモリを待つスレッドはスリープしているのではなく，
オペレーティングシステムからは見えないプロセッサの内部で待っているから
である．個々のメモリアクセスをソフトウェアで追跡することもできない．それ
には時間がかかりすぎる．情報はハードウェアから得なければならない．

\begin{definition}[ハードウェア性能カウンタ]
  \term[はーどうぇあせいのうかうんた]{ハードウェア性能カウンタ}%
  [hardware performance counter]とは，プロセッサの実行に伴って更新され，
  性能に関係する事象を数えるプロセッサのレジスタである．
\end{definition}

プロセッサは通常いくつかのカウンタを備えており，利用できるカウンタは
アーキテクチャによって異なる．例えば次のようなものを数える．
\begin{itemize}
  \item L1，L2 から LLC に至るキャッシュ階層の各水準でのキャッシュミス．
  \item 命令数とサイクル数．ここからサイクル当たり命令数（IPC）や，その
    逆数である命令当たりサイクル数が導かれる．
  \item 一部のプロセッサでは電力とエネルギーの消費．
\end{itemize}
ソフトウェアのインタフェースやツールがこれらのカウンタを利用可能にする．
例えば，アーキテクチャごとに利用できるカウンタを文書に列挙している
\texttt{oprofile} や，Linux の \texttt{perf} ツールがある．\tip{コマンド
\texttt{perf stat -e cycles,instructions,LLC-load-misses} \emph{cmd} は，
プログラム \emph{cmd} についてこれらの計数を報告する．}

カウンタの値から，スケジューラはタスクがどのような資源を必要とするかを
推定できる．LLC ミスの少ないタスクはデータをキャッシュ内に見出しており
CPU バウンドである可能性が高く，LLC ミスの多いタスクはメモリを待ち続けて
おりメモリバウンドである可能性が高い．単一のカウンタで決着が付くことは
まれである．そこでスケジューラは複数のカウンタを組み合わせ，よく理解された
ワークロードの特性に基づいてアーキテクチャごとに閾値を調整したモデルを
用いる．

\begin{definition}[命令当たりサイクル数]
  タスクの \term{CPI}[cycles per instruction]%
  \index{cycles per instruction|see{CPI}}（命令当たりサイクル数）とは，
  実行した命令1つ当たりに要する CPU サイクル数，すなわちサイクル数の計数
  と命令数の計数の比である．
\end{definition}

\begin{example}
  \SI{100}{\milli\second} の区間の間に，カウンタが，タスク A は
  $4.6 \times 10^8$ サイクルで $4.2 \times 10^8$ 命令を実行し，1000 命令
  当たりの LLC ミスは 0.2 であったと報告し，タスク B は同じサイクル数で
  $0.9 \times 10^8$ 命令を実行し，1000 命令当たりの LLC ミスは 25 であった
  と報告したとする．A の CPI は約 1.1，B の CPI は約 5.1 である．どちらの
  指標も A を CPU バウンド，B をメモリバウンドと分類しており，
  \cref{sec:l07-smt-sched}に従うスケジューラなら，これらを同じ CPU の
  2つのハードウェアスレッド上に配置するであろう．
\end{example}

\needspace{7\baselineskip}
\section{命令当たりサイクル数は有用な指標か}
\label{sec:l07-cpi}

Fedorova らは，カウンタから導かれる単一の指標である CPI によってタスクを
区別することを提案した．\source{P3L1，CPI の実験，CPI の実験結果；
Fedorova ら．}この着想は魅力的である．メモリバウンドのタスクは多くの
サイクルをメモリ待ちに費やすので CPI が高く，CPU バウンドのタスクは CPI
が低い．1サイクルに1命令を実行するモデルでは 1 であり，スーパースカラの
プロセッサでは 1 を下回る．そうであれば，スケジューラは各コアの
ハードウェアスレッド上に CPI の異なるタスクを混ぜればよいことになる．

この着想を検証するために，彼らは4つのコアがそれぞれ4つのハードウェア
スレッドを持ち，合計 16 のハードウェアスレッドを持つチップマルチ
スレッディングのプロセッサをシミュレートした．ワークロードは合成されたもの
で，CPI が大きく異なるスレッドからなっていた．スレッドは，どのコアも
さまざまな CPI のスレッドを混ぜて実行する割り当てから，どのコアも同じ CPI
のスレッドを実行する割り当てまで，いくつかの方法でコアへ割り当てられた．
性能は，16 のハードウェアスレッドの合計のサイクル当たり命令数として測られ，
4つのコアでは最大 4 である．

CPI を混ぜた割り当てが最も高い IPC を達成し，各コア上のスレッドの CPI が
似ている割り当てが最も低い IPC となった．これらのワークロードでは CPI に
よるスケジューリングが明らかに有益であり，この実験はプロセッサの
パイプラインの奪い合いが重要であることを裏付けている．しかし現実的な
ワークロードは異なる振る舞いを示す．標準的なベンチマークスイートのような
実際のアプリケーションについて測定された CPI は，それほど明確には分かれ
ない．それらは狭い範囲に集中する．\sidenote{ワークショップ論文自体は CPI
について楽観的であった．測定された CPI はアプリケーション間で1桁の幅で
変動し，スレッドの近い将来の CPI は予測可能であると報告している．そこで
挙げられている限界は別のものである．CPI はキャッシュの奪い合いも浮動小数点
ユニットの奪い合いも明らかにせず，忙しいシステム上で測ると歪む．}
したがって CPI では，現実のワークロードのタスクを，スケジューラを導くのに
十分なほど区別することはできない．合成されたワークロードでうまくいくことが，
実際にはうまくいかないのである．

\begin{note}
  それでもこの研究は3つの教訓を残している．SMT のプロセッサ上のタスクは
  プロセッサのパイプラインを奪い合うこと，ハードウェアカウンタは
  ワークロードを特徴付けられること，そしてスケジューラは単に負荷を均衡
  させるだけでなく，資源の奪い合いを意識しなければならないことである．
  同じグループによる後の研究は，奪い合いの指標としては CPI より，ラスト
  レベルキャッシュの時間的な使用状況の方が優れていることを見出した．
\end{note}

\begin{keypoint}[章のまとめ]
  スケジューラは自分のランキューから次のタスクを選び，ランキューの構造は
  アルゴリズムとともに設計される．FCFS と SJF はタスクを完了まで実行する．
  プリエンプションと優先度は緊急なタスクに応えるが，エージングで回避
  されるスタベーションと，優先度継承で回避される優先度逆転をもたらす．
  タイムスライスを用いるラウンドロビンはシステムを応答性の高いものにする．
  CPU バウンドのタスクは長いタイムスライスを，I/O バウンドのタスクは短い
  タイムスライスを好み，多段フィードバックキューは各タスクの種類を学習
  することでそれを実現する．Linux は，アクティブ配列と期限切れ配列を持つ
  $O(1)$ スケジューラを，赤黒木から仮想実行時間が最小のタスクを実行する
  CFS に置き換えた．マルチプロセッサでは，CPU ごとのランキューがキャッシュ
  アフィニティを保ち，負荷分散がタスクを分散させ，NUMA 対応スケジューリング
  がタスクをそのメモリの近くに保つ．SMT のプロセッサでは，CPU バウンドの
  タスクとメモリバウンドのタスクが CPU を共有すべきであり，ハードウェア
  性能カウンタがそれらを区別する助けとなるが，CPI だけでは十分な指標では
  ない．
\end{keypoint}

\end{document}
